\cleardoublepage

\oldpage[IV-1]{222}
\thispagestyle{plain}

\begin{center}
{\Large CHAPTER IV}

\vspace{1.5cm}
{\LARGE\bfseries LOCAL STUDY OF SCHEMES\\[0.3em]
AND OF MORPHISMS OF SCHEMES}

\vspace{1.5cm}
{\large\bfseries Summary}
\end{center}

\medskip
\begin{tabular}{@{}r p{0.84\textwidth}}
§ 1. & Relative finiteness conditions. Constructible sets in preschemes.\\
§ 2. & Base change and flatness.\\
§ 3. & Associated prime cycles and primary decompositions.\\
§ 4. & Change of the base field in preschemes.\\
§ 5. & Dimension and depth in preschemes.\\
§ 6. & Flat morphisms of locally Noetherian preschemes.\\
§ 7. & Application to the relations between a Noetherian local ring and its completion.
Excellent rings.\\
§ 8. & Projective limits of preschemes.\\
§ 9. & Constructible properties.\\
§ 10. & Jacobson preschemes.\\
§ 11.\footnotemark[1] & Topological properties of flat morphisms of finite presentation.
Local flatness criteria.\\
§ 12. & Study of the fibres of flat morphisms of finite presentation.\\
§ 13. & Equidimensional morphisms.\\
§ 14. & Universally open morphisms.\\
§ 15. & Study of the fibres of a universally open morphism.\\
§ 16. & Differential invariants. Differentially smooth morphisms.\\
§ 17. & Smooth, unramified, étale morphisms.\\
§ 18. & Supplement on étale morphisms. Henselian local rings.\\
§ 19. & Regular and transversally regular immersions.\\
§ 20. & Hyperplane sections; generic projections.\\
§ 21. & Infinitesimal extensions.\\
\end{tabular}
\footnotetext[1]{The order and content of §§ 11 to 21 are given only as an
indication and may be modified before their publication.}

\clearpage
\oldpage[IV-1]{223}

The subjects treated in this chapter call for the following remarks:

\medskip
\noindent
\emph{a)} Their common feature is that they concern \emph{local} properties
of preschemes or of morphisms, i.e. properties considered at a point, or at
the points of a fibre, or in an (unspecified) neighbourhood of a point or of
a fibre. These properties are generally of a \emph{topological},
\emph{differential} or \emph{dimensional} nature (i.e. involving the
notions of \emph{dimension} and of \emph{depth}) and are linked to the
properties of the \emph{local rings} at the points under consideration. A
typical problem is to relate, for a given morphism $f:X\to Y$ and a point
$x\in X$, the properties of $X$ at $x$ to those of $Y$ at $y=f(x)$ and of
the fibre $X_y=f^{-1}(y)$ at $x$. Another is to determine the topological
nature (for example constructibility, or the fact of being open or
closed) of the set of points $x\in X$ at which $X$ has a certain
property or for which the fibre $X_{f(x)}$ passing through $x$ has a
certain property at $x$. Likewise, we will be interested in the topological
nature of the set of points $y\in Y$ such that $X$ has a certain property at
all the points of the fibre $X_y$, or such that this fibre has a certain
property.

\medskip
\noindent
\emph{b)} The most important notions for the following chapters are
that of a \emph{flat morphism of finite presentation} and those of a
\emph{smooth morphism} and of an \emph{étale morphism}, which are particular
cases of it. Their detailed study (together with that of related questions)
begins, properly speaking, in §~11.

\medskip
\noindent
\emph{c)} §§~1 to 10 may be regarded as preliminary in nature
and develop three types of techniques, used not only in the
other sections of the chapter but also, of course, in the following
chapters:

\medskip
\noindent
\emph{c 1)} In §§~1 to 4, various aspects of the notion of
\emph{base change} are considered, above all in relation to
\emph{finiteness} or \emph{flatness} conditions; there we begin the
technique of \emph{descent} in its most elementary aspects (the questions
of ``effectiveness'' linked to this technique will be studied in chap.~V).

\medskip
\noindent
\emph{c 2)} §§~5 to 7 are centred on what may be called
\emph{Noetherian techniques}, the preschemes considered always being
supposed locally Noetherian, whereas, on the contrary, no finiteness
condition is generally imposed on the \emph{morphisms}; this is
essentially because the notions of dimension and depth are hardly
manageable except for Noetherian local rings. Let us recall that §~7
constitutes a ``fine'' theory of Noetherian local rings, rather little
used in the rest of the chapter.

\medskip
\noindent
\emph{c 3)} §§~8 to 10 give, among other things, the means of
\emph{eliminating Noetherian hypotheses} on the preschemes studied, by
substituting for them suitable \emph{finiteness} hypotheses (``finite
presentation'') on the \emph{morphisms} considered: the advantage of this
substitution is that the latter hypotheses are \emph{stable under base
change}, which is not the case for Noetherian hypotheses on preschemes. The
technique that makes this passage possible rests, on the one hand, on the use
of the notion of \emph{projective limit} of preschemes, thanks to which one
can reduce a question to the same question under \emph{Noetherian}
hypotheses; on the other hand, on the systematic use of
\emph{constructible sets}, which have the double advantage of being preserved
under inverse image (by an arbitrary

\oldpage[IV-1]{224}

morphism) and under direct image (by a morphism of finite presentation), and
of having manageable topological properties in locally Noetherian
preschemes. The same techniques often even make it possible to reduce to
more special Noetherian rings, for example the
$\mathbf Z$-\emph{algebras of finite type}, and it is there that the
properties of ``excellent'' rings (studied in §~7) play a decisive role.
Moreover, independently of the question of eliminating Noetherian
hypotheses, the techniques of §§~8 to 10, which are of a very elementary
nature, are in constant use in almost all applications.

\setcounter{section}{0}


\section{Relative finiteness conditions. Constructible sets in preschemes}
\label{section:IV.1}

In this section we take up again, and complete, the exposition of the
``finiteness conditions'' for a morphism of preschemes $f:X\to Y$ given
in (I, 6.3 and 6.6). There are essentially two notions of ``finiteness'' of
a \emph{global} nature on $X$, that of a \emph{quasi-compact} morphism
(defined in (I, 6.6.1)) and that of a \emph{quasi-separated} morphism; on
the other hand, there are two notions of ``finiteness'' of a \emph{local}
nature on $X$, that of a morphism \emph{locally of finite type} (defined in
(I, 6.6.2)) and that of a morphism \emph{locally of finite presentation}. By
combining these local notions with the preceding global notions, we obtain
the notions of a morphism \emph{of finite type} (defined in (I, 6.3.1)) and
of a morphism \emph{of finite presentation}. For the reader's convenience,
we will give again in this section the main properties stated in (I, 6.3
and 6.6), referring of course to those subsections of Chapter I for their
proofs.

In n\textsuperscript{os}~(1.8) and (1.9) we complete, in the setting of
preschemes and making use of the preceding finiteness notions, the
results on constructible sets given in
(0\textsubscript{III}, §~9).

\subsection{Quasi-compact morphisms}
\label{subsection:IV.1.1}

\begin{definition}[1.1.1]
\label{IV.1.1.1}
We say that a morphism of preschemes $f:X\to Y$ is \emph{quasi-compact} if
the continuous map $f$ from the topological space $X$ to the topological
space $Y$ is quasi-compact (0\textsubscript{III}, 9.1.1), in other
words if the inverse image $f^{-1}(U)$ of every quasi-compact open subset $U$
of $Y$ is quasi-compact (cf. (I, 6.6.1)).
\end{definition}

If $\mathfrak B$ is a basis of the topology of $Y$ consisting of affine open
subsets, then for $f$ to be quasi-compact it is necessary and sufficient that
for every $V\in\mathfrak B$, $f^{-1}(V)$ be a \emph{finite union of affine
open subsets}. For example, if $Y$ is affine and $X$ quasi-compact,
\emph{every} morphism $f:X\to Y$ is quasi-compact (I, 6.6.1).

If $f:X\to Y$ is a quasi-compact morphism, it is clear that for every
open subset $V$ of $Y$, the restriction of $f$ to $f^{-1}(V)$ is a
quasi-compact morphism $f^{-1}(V)\to V$. Conversely, if $(U_\alpha)$ is an
open covering of $Y$ and $f:X\to Y$ a morphism such that the restrictions
$f^{-1}(U_\alpha)\to U_\alpha$ are quasi-compact, then $f$ is
quasi-compact. Consequently,

\oldpage[IV-1]{225}

if $f:X\to Y$ is an $S$-morphism of $S$-preschemes, and if there exists
an open covering $(S_\lambda)$ of $S$ such that the restrictions
$g^{-1}(S_\lambda)\to h^{-1}(S_\lambda)$ of $f$ (where $g$ and $h$ are the
structural morphisms) are quasi-compact, then $f$ is quasi-compact.

\begin{proposition}[1.1.2]
\label{IV.1.1.2}
\begin{enumerate}
  \item[(i)] An immersion $X\to Y$ is quasi-compact if it is closed, or if
  the underlying space of $Y$ is locally Noetherian, or if the underlying
  space of $X$ is Noetherian.
  \item[(ii)] The composite of two quasi-compact morphisms is quasi-compact.
  \item[(iii)] If $f:X\to Y$ is a quasi-compact $S$-morphism, then so is
  $f_{(S')}:X_{(S')}\to Y_{(S')}$ for every extension $g:S'\to S$ of the
  base prescheme.
  \item[(iv)] If $f:X\to X'$ and $g:Y\to Y'$ are two quasi-compact
  $S$-morphisms,
  \[
    f\times_S g:X\times_S Y\to X'\times_S Y'
  \]
  is quasi-compact.
  \item[(v)] If the composite $g\circ f$ of two morphisms $f:X\to Y$, $g:Y\to Z$
  is quasi-compact, and if $g$ is separated, or the underlying space of $X$
  locally Noetherian, then $f$ is quasi-compact.
  \item[(vi)] For a morphism $f$ to be quasi-compact, it is necessary and
  sufficient that $f_{\mathrm{red}}$ be so.
\end{enumerate}
\end{proposition}

For the proof, see (I, 6.6.4). Note that assertion (vi) is
also a consequence of the more general proposition:

\begin{proposition}[1.1.3]
\label{IV.1.1.3}
Let $f:X\to Y$, $g:Y\to Z$ be two morphisms. If $g\circ f$ is
quasi-compact and $f$ surjective, then $g$ is quasi-compact.
\end{proposition}

\begin{proof}
Indeed, if $V$ is a quasi-compact open subset of $Z$,
$f^{-1}(g^{-1}(V))$ is quasi-compact by hypothesis, and we have
$g^{-1}(V)=f(f^{-1}(g^{-1}(V)))$, since $f$ is surjective; hence
$g^{-1}(V)$ is quasi-compact.
\end{proof}

\begin{corollary}[1.1.4]
\label{IV.1.1.4}
Let $f:X\to Y$, $g:Y'\to Y$ be two morphisms, $X'=X\times_Y Y'$,
$f'=f_{(Y')}:X'\to Y'$; suppose $g$ quasi-compact and surjective (resp.
surjective). Then, for $f$ to be quasi-compact (resp. dominant), it suffices
that $f'$ be so.
\end{corollary}

\begin{proof}
If $g':X'\to X$ is the canonical projection, $g'$ is a surjective morphism
(I, 3.5.2), and we have $f\circ g'=g\circ f'$; if $f'$ is quasi-compact
(resp. dominant), then so is $g\circ f'$, since $g$ is
quasi-compact (1.1.2) (resp. surjective); since $f\circ g'$ is quasi-compact
(resp. dominant) and $g'$ is surjective, we deduce that $f$ is
quasi-compact (1.1.3) (resp. dominant).
\end{proof}

\oldpage[IV-1]{226}

For brevity, we will call a \emph{maximal point} of a prescheme $X$ any
generic point of an irreducible component of $X$; if
$X=\operatorname{Spec}(A)$ is affine, the maximal points of $X$ are the
\emph{minimal prime ideals} of $A$.

\begin{proposition}[1.1.5]
\label{IV.1.1.5}
Let $f:X\to Y$ be a quasi-compact morphism. The following conditions are
equivalent:
\begin{enumerate}
  \item[a)] $f$ is dominant;
  \item[b)] for every maximal point $y$ of $Y$, we have
  $f^{-1}(y)\neq\varnothing$;
  \item[c)] for every maximal point $y$ of $Y$, $f^{-1}(y)$ contains a
  maximal point of $X$.
\end{enumerate}
\end{proposition}

The equivalence of \emph{a)} and \emph{c)} was proved in (I, 6.6.5). It
is clear that \emph{c)} implies \emph{b)}; on the other hand, if $X'$ is an
irreducible component of $X$ meeting $f^{-1}(y)$, $f(X')$ is
irreducible and contains $y$, hence is contained in the irreducible component
$Y'=\overline{\{y\}}$ of $Y$ (0\textsubscript{I}, 2.1.6); if $x$ is the
generic point of $X'$, we therefore necessarily have $f(x)=y$
(0\textsubscript{I}, 2.1.5); hence \emph{b)} implies \emph{c)}.

\begin{proposition}[1.1.6]
\label{IV.1.1.6}
Let $f':X'\to Y$, $f'':X''\to Y$ be two morphisms, $X$ the sum prescheme
$X'\coprod X''$, $f:X\to Y$ the morphism equal to $f'$ on $X'$ and to $f''$
on $X''$. For $f$ to be quasi-compact, it is necessary and sufficient that
$f'$ and $f''$ be so.
\end{proposition}

This follows immediately from the definition.

\subsection{Quasi-separated morphisms}
\label{subsection:IV.1.2}

\begin{definition}[1.2.1]
\label{IV.1.2.1}
Let $X$, $Y$ be two preschemes. We say that a morphism $f:X\to Y$ is
\emph{quasi-separated} (or that $X$ is \emph{quasi-separated over} $Y$) if
the diagonal morphism
$\Delta_f=(1_X,1_X)_Y:X\to X\times_Y X$ is quasi-compact. We say that a
prescheme $X$ is \emph{quasi-separated} if it is quasi-separated over
$\operatorname{Spec}(\mathbf Z)$.
\end{definition}

By definition, every separated morphism is quasi-separated, $\Delta_f$ being a
closed immersion (1.1.2, (i)); a \emph{scheme} $X$ is a quasi-separated
prescheme, being separated over $\operatorname{Spec}(\mathbf Z)$ (I, 5.4.1).

\begin{proposition}[1.2.2]
\label{IV.1.2.2}
\begin{enumerate}
  \item[(i)] Every monomorphism of preschemes $f:X\to Y$ (in particular every
  immersion) is quasi-separated.
  \item[(ii)] If $f:X\to Y$ and $g:Y\to Z$ are two quasi-separated morphisms,
  $g\circ f:X\to Z$ is quasi-separated.
  \item[(iii)] Let $X$, $Y$ be two $S$-preschemes, $f:X\to Y$ a quasi-separated
  $S$-morphism. Then, for every base change $g:S'\to S$, the morphism
  $f_{(S')}:X_{(S')}\to Y_{(S')}$ is quasi-separated.
  \item[(iv)] If $f:X\to Y$, $f':X'\to Y'$ are two quasi-separated $S$-morphisms,
  then
  \[
    f\times_S f':X\times_S X'\to Y\times_S Y'
  \]
  is quasi-separated.
  \item[(v)] If $f:X\to Y$, $g:Y\to Z$ are two morphisms such that $g\circ f$
  is quasi-separated, then $f$ is quasi-separated.
  \item[(vi)] If $f:X\to Y$ is a quasi-separated morphism,
  $f_{\mathrm{red}}:X_{\mathrm{red}}\to Y_{\mathrm{red}}$ is quasi-separated.
\end{enumerate}
\end{proposition}

\begin{proof}
By virtue of (I, 5.5.12), it suffices to prove (i), (ii) and (iii).

Assertion (i) is trivial, every monomorphism being separated (I, 5.5.1).
To prove (iii), we immediately reduce to the case where $Y=S$ (I, 3.3.11) and
the assertion follows from the fact that
$\Delta_{f_{(S')}}=(\Delta_f)_{(S')}$ (I, 5.3.4) and from (1.1.2, (iii)). To
prove (ii), consider the projections $p$ and $q$ of $X\times_Y X$ onto
$X$; if $i=(p,q)_Z$, we know that the diagram
\[
\xymatrix{
  X\times_Y X \ar[r]^-{i} \ar[d]_-{\pi} & X\times_Z X
    \ar[d]^-{f\times_Z f} \\
  Y \ar[r]_-{\Delta_g} & Y\times_Z Y
}
\]
\oldpage[IV-1]{227}
(where $\pi$ is the structural morphism) is commutative and identifies
$X\times_Y X$ with the product of the $(Y\times_Z Y)$-preschemes $Y$ and
$X\times_Z X$ (I, 5.3.5). If $g$ is quasi-separated, $\Delta_g$ is
quasi-compact, hence so is $i$ (1.1.2, (iii)); if moreover $f$
is quasi-separated, $\Delta_f$ is quasi-compact, hence so is
$i\circ\Delta_f$ (1.1.2, (ii)), which is equal to $\Delta_{g\circ f}$.
\end{proof}

\begin{corollary}[1.2.3]
\label{IV.1.2.3}
\begin{enumerate}
  \item[(i)] Let $X$ be a quasi-separated prescheme. Then every morphism
  $f:X\to Y$ is quasi-separated.
  \item[(ii)] Let $Y$ be a quasi-separated prescheme. For a morphism
  $f:X\to Y$ to be quasi-separated, it is necessary and sufficient that the
  prescheme $X$ be quasi-separated.
\end{enumerate}
\end{corollary}

This follows immediately from (1.2.2, (ii) and (v)) applied to $X$, $Y$ and
$\operatorname{Spec}(\mathbf Z)$.

\begin{proposition}[1.2.4]
\label{IV.1.2.4}
Let $f:X\to Y$ be a morphism, $g:Y\to Z$ a quasi-separated morphism. If
$g\circ f$ is quasi-compact, then so is $f$.
\end{proposition}

\begin{proof}
We know that $f$ is the composite morphism
$X\xrightarrow{\Gamma_f}X\times_ZY\xrightarrow{\mathrm{pr}_2}Y$
(I, 5.3.13); on the other hand, $\mathrm{pr}_2$ is identified with
$(g\circ f)\times_Z1_Y$ (I, 3.3.4), and if $g\circ f$ is quasi-compact, so
is $\mathrm{pr}_2$ (1.1.2, (iv)). Finally, we have the commutative
diagram
\[
\xymatrix{
  X \ar[r]^-{\Gamma_f} \ar[d]_-{f} & X\times_ZY
    \ar[d]^-{f\times 1_Y} \\
  Y \ar[r]_-{\Delta_g} & Y\times_ZY
}
\tag{1.2.4.1}
\label{IV.1.2.4.1}
\]
which identifies $X$ with the product of the $(Y\times_ZY)$-preschemes $Y$ and
$X\times_ZY$ (I, 5.3.7). Since by hypothesis $\Delta_g$ is a quasi-compact
morphism, so is $\Gamma_f$ (1.1.2, (iii)); the conclusion
therefore follows from (1.1.2, (ii)).
\end{proof}

\begin{proposition}[1.2.5]
\label{IV.1.2.5}
Let $f:X\to Y$, $g:Y'\to Y$ be two morphisms,
$X'=X\times_Y Y'$, $f'=f_{(Y')}:X'\to Y'$. If $g$ is quasi-compact and
surjective and $f'$ quasi-separated, then $f$ is quasi-separated.
\end{proposition}

\begin{proof}
Indeed, we have $(X'\times_{Y'}X')=(X\times_YX)_{(Y')}$ and
$\Delta_{f'}=(\Delta_f)_{(Y')}$ (I, 5.3.4); since the projection
$X'\times_{Y'}X'\to X\times_YX$ is quasi-compact (1.1.2, (iii)) and
surjective (I, 3.5.2), we may, by virtue of (I, 3.3.11), apply (1.1.4),
which shows that since $\Delta_{f'}$ is a quasi-compact morphism, so is
$\Delta_f$.
\end{proof}

\begin{proposition}[1.2.6]
\label{IV.1.2.6}
Let $f:X\to Y$ be a morphism, $(U_\alpha)$ a covering of $Y$ by open
subsets such that the induced preschemes $U_\alpha$ are quasi-separated. For
$f$ to be quasi-separated, it is necessary and sufficient that each of the
preschemes induced by $X$ on the open subsets $f^{-1}(U_\alpha)$ be
quasi-separated.
\end{proposition}

\begin{proof}
The inverse image of $U_\alpha$ in $X\times_YX$ is
$X_\alpha\times_{U_\alpha}X_\alpha$, where we set
$X_\alpha=f^{-1}(U_\alpha)$, and the restriction
$X_\alpha\to X_\alpha\times_{U_\alpha}X_\alpha$ of $\Delta_f$ is none other
than $\Delta_{f_\alpha}$, denoting by $f_\alpha$ the restriction
$X_\alpha\to U_\alpha$ of $f$; by virtue of Definition (1.2.1) and of the
local character of the notion of quasi-compact morphism (1.1.1), for $f$
to be quasi-separated, it is necessary and sufficient that each of the
morphisms $f_\alpha$ be so. But since by hypothesis the morphism
$U_\alpha\to\operatorname{Spec}(\mathbf Z)$ is quasi-separated, it amounts to
the same to say that $f_\alpha$ is quasi-separated or that the composite
$X_\alpha\xrightarrow{f_\alpha}U_\alpha\to\operatorname{Spec}(\mathbf Z)$
is so (1.2.2, (ii) and (v)), whence the proposition.
\end{proof}

\oldpage[IV-1]{228}

To verify that a morphism is quasi-separated, we are therefore reduced to
giving criteria for a prescheme to be quasi-separated:

\begin{proposition}[1.2.7]
\label{IV.1.2.7}
Let $X$ be a prescheme, $(U_\alpha)$ a covering of $X$ consisting of
quasi-compact open subsets. The following properties are equivalent:
\begin{enumerate}
  \item[a)] $X$ is quasi-separated.
  \item[b)] For every quasi-compact open subset $U$ of $X$, the canonical
  immersion $U\to X$ is quasi-compact (in other words, $U$ is retrocompact
  in $X$).
  \item[b')] The intersection of two quasi-compact open subsets of $X$ is
  quasi-compact.
  \item[c)] For every pair of indices $(\alpha,\beta)$, the intersection
  $U_\alpha\cap U_\beta$ is quasi-compact.
\end{enumerate}
\end{proposition}

\begin{proof}
Properties \emph{b)} and \emph{b')} are trivially equivalent by the
definition of a quasi-compact morphism. Since a quasi-compact open subset is a
finite union of affine open subsets, for two quasi-compact open subsets $U$,
$V$ of $X$, $U\times_{\mathbf Z}V$ is a quasi-compact open subset of
$X\times_{\mathbf Z}X$ (I, 3.2.7), whose inverse image under $\Delta_X$ is
$U\cap V$; hence \emph{a)} implies \emph{b')}. It is trivial that \emph{b')}
implies \emph{c)}; finally, if \emph{c)} holds, the
$U_\alpha\times_{\mathbf Z}U_\beta$ form a covering of
$X\times_{\mathbf Z}X$ by quasi-compact open subsets, and we know that for
the morphism $\Delta_X:X\to X\times_{\mathbf Z}X$ to be quasi-compact, it
suffices that the inverse images $U_\alpha\cap U_\beta$ of the
$U_\alpha\times_{\mathbf Z}U_\beta$ under $\Delta_X$ be quasi-compact
(1.1.1); hence \emph{c)} implies \emph{a)}.
\end{proof}

\begin{corollary}[1.2.8]
\label{IV.1.2.8}
Every prescheme $X$ whose underlying space is locally Noetherian (for
example a locally Noetherian prescheme) is quasi-separated; every morphism
$f:X\to Y$ is then quasi-separated.
\end{corollary}

\begin{proposition}[1.2.9]
\label{IV.1.2.9}
Let $X'$, $X''$ be two preschemes, $X=X'\coprod X''$ their sum,
$f':X'\to Y$, $f'':X''\to Y$ two morphisms, $f:X\to Y$ the morphism that
coincides with $f'$ on $X'$ and with $f''$ on $X''$. For $f$ to be
quasi-separated it is necessary and sufficient that $f'$ and $f''$ be so.
\end{proposition}

\begin{proof}
Indeed, $X\times_YX$ is the sum of the four preschemes
$X'\times_YX'$, $X'\times_YX''$, $X''\times_YX'$ and $X''\times_YX''$, and
$\Delta_f$ is the morphism that coincides with $\Delta_{f'}$ on $X'$ and with
$\Delta_{f''}$ on $X''$; the proposition therefore follows immediately from
the definitions.
\end{proof}

\subsection{Morphisms locally of finite type}
\label{subsection:IV.1.3}

\begin{env}[1.3.1]
\label{IV.1.3.1}
Recall that if $A$ is a ring, a (commutative) $A$-algebra $B$ is said to be
\emph{of finite type} if it is generated by finitely many elements of
$B$, or, what amounts to the same, if it is isomorphic to the quotient of a
polynomial algebra $A[T_1,\ldots,T_n]$ by an ideal of this algebra. It is
clear that for every (commutative) $A$-algebra $A'$,
$B\otimes_AA'$ is then an $A'$-algebra of finite type. If $B$ is an
$A$-algebra of finite type and $C$ a $B$-algebra of finite type, $C$ is an
$A$-algebra of finite type, for if $C$ is a quotient of
$B[T_1,\ldots,T_n]$ and $B$ a quotient of $A[S_1,\ldots,S_m]$,
$B[T_1,\ldots,T_n]=B\otimes_AA[T_1,\ldots,T_n]$ is a quotient of
$A[S_1,\ldots,S_m,T_1,\ldots,T_n]$, hence so is $C$.
\end{env}

\begin{definition}[1.3.2]
\label{IV.1.3.2}
Let $f:X\to Y$ be a morphism of preschemes, $x$ a point of $X$,
$y=f(x)$. We say that $f$ is \emph{of finite type at $x$} if there exist an
affine open neighbourhood $V$ of $y$ and an affine open neighbourhood $U$ of
$x$ such that $f(U)\subset V$ and the ring $A(U)$ is an $A(V)$-algebra of
finite type. We say that $f$ is \emph{locally of finite type} if it is of
finite type at every point of $X$.
\end{definition}

\oldpage[IV-1]{229}

\begin{proposition}[1.3.3]
\label{IV.1.3.3}
If $Y$ is a locally Noetherian prescheme and $f:X\to Y$ a morphism
locally of finite type, then $X$ is locally Noetherian.
\end{proposition}

For the proof, see (I, 6.3.7).

\begin{proposition}[1.3.4]
\label{IV.1.3.4}
\begin{enumerate}
  \item[(i)] Every local immersion is locally of finite type.
  \item[(ii)] If two morphisms $f:X\to Y$, $g:Y\to Z$ are locally of
  finite type, then so is $g\circ f$.
  \item[(iii)] If $f:X\to Y$ is an $S$-morphism locally of finite type,
  $f_{(S')}:X_{(S')}\to Y_{(S')}$ is locally of finite type for every
  extension $S'\to S$ of the base prescheme.
  \item[(iv)] If $f:X\to X'$ and $g:Y\to Y'$ are two $S$-morphisms
  locally of finite type, $f\times_Sg$ is locally of finite type.
  \item[(v)] If the composite $g\circ f$ of two morphisms is locally of
  finite type, $f$ is locally of finite type.
  \item[(vi)] If a morphism $f$ is locally of finite type, then so is
  $f_{\mathrm{red}}$.
\end{enumerate}
\end{proposition}

For the proof, see (I, 6.6.6).

\begin{corollary}[1.3.5]
\label{IV.1.3.5}
Let $f:X\to Y$ be a morphism locally of finite type. For every morphism
$Y'\to Y$ such that $Y'$ is locally Noetherian, $X\times_YY'$ is
locally Noetherian.
\end{corollary}

\begin{proof}
Indeed, $f_{(Y')}:X\times_YY'\to Y'$ is locally of finite type by
(1.3.4, (iii)), and it suffices to apply (1.3.3).
\end{proof}

\begin{proposition}[1.3.6]
\label{IV.1.3.6}
Let $\varphi:A\to B$ be a ring homomorphism. For the corresponding
morphism $f:\operatorname{Spec}(B)\to\operatorname{Spec}(A)$ to be locally of finite type, it is necessary and
sufficient that $B$ be an $A$-algebra of finite type.
\end{proposition}

For the proof, see (I, 6.3.3), taking into account the fact that $f$ is
quasi-compact, and (I, 6.6.3).

\begin{proposition}[1.3.7]
\label{IV.1.3.7}
For a morphism locally of finite type $f:X\to Y$ to be surjective, it is
necessary and sufficient that, for every algebraically closed field $\Omega$,
the map $X(\Omega)\to Y(\Omega)$ corresponding to $f$ (I, 3.4.1) be
surjective.
\end{proposition}

For the proof, see (I, 6.3.10).

\begin{env}[1.3.8]
\label{IV.1.3.8}
Let $A$ be a ring. We say that an $A$-algebra $B$ is \emph{essentially of
finite type} if $B$ is $A$-isomorphic to an $A$-algebra of the form $S^{-1}C$,
where $C$ is an $A$-algebra \emph{of finite type} and $S$ a multiplicative
subset of $C$.
\end{env}

\par\medskip\noindent
\begin{proposition}[1.3.9]
\label{IV.1.3.9}
\begin{enumerate}
  \item[(i)] If $B$ is an $A$-algebra essentially of finite type and $C$
  a $B$-algebra essentially of finite type, then $C$ is an
  $A$-algebra essentially of finite type.
  \item[(ii)] Let $B$ be an $A$-algebra essentially of finite type, $A'$
  an $A$-algebra; then $B'=B\otimes_AA'$ is an $A'$-algebra
  essentially of finite type.
\end{enumerate}
\end{proposition}

\begin{proof}
\begin{enumerate}
  \item[(i)] Let $B=S'^{-1}B'$ and $C=T'^{-1}C'$, where $B'$ (resp. $C'$) is
  an $A$-algebra (resp. a $B$-algebra) of finite type and $S'$ (resp. $T'$)
  a multiplicative subset of $B'$ (resp. $C'$); then $C'$ is of the form
  $S'^{-1}C''$, where $C''$ is a $B'$-algebra of finite type, hence $C$ is
  also a ring of fractions of $C''$; since $C''$ is an $A$-algebra of finite
  type, this proves our assertion.
  \item[(ii)] If $B$ is a ring of fractions of an $A$-algebra of finite type
  $C$, $B'$ is a ring of fractions of the $A'$-algebra of finite type
  $C\otimes_AA'$, whence (ii).
\end{enumerate}
\end{proof}

\oldpage[IV-1]{230}

\begin{env}[1.3.10]
\label{IV.1.3.10}
If $B$ is a \emph{local} $A$-algebra essentially of finite type, $B$ is
of the form $C_{\mathfrak q}$, where $C$ is an $A$-algebra of finite type and
$\mathfrak q$ a prime ideal of $C$ (Bourbaki, \emph{Alg. comm.}, chap. II,
\S~5, prop. 11). Let $\mathfrak p$ be the inverse image in $A$ of the ideal
$\mathfrak q$; if we set $S=A-\mathfrak p$, $C_{\mathfrak q}$ is also a
local ring at a prime ideal of $S^{-1}C$; since $S^{-1}C$ is an
algebra of finite type over $A_{\mathfrak p}=S^{-1}A$, we see that $B$ is also
an $A_{\mathfrak p}$-algebra \emph{essentially of finite type}, and moreover
the homomorphism $A_{\mathfrak p}\to B$ is \emph{local}.
\end{env}

\begin{proposition}[1.3.11]
\label{IV.1.3.11}
If $B$ is a local $A$-algebra essentially of finite type, $B$ is
$A$-isomorphic to a quotient $A$-algebra of an $A$-algebra of the form
$C_{\mathfrak q}$, where $C$ is a polynomial algebra
$A[T_1,\ldots,T_n]$ and $\mathfrak q$ a prime ideal of $C$.
\end{proposition}

\begin{proof}
Indeed, by definition, $B$ is isomorphic to $C'_{\mathfrak q'}$, where $C'$ is
an $A$-algebra of finite type and $\mathfrak q'$ a prime ideal of $C'$; but
$C'=C/\mathfrak b$, where $C=A[T_1,\ldots,T_n]$ and $\mathfrak b$ is an ideal
of $C$; hence $\mathfrak q'=\mathfrak q/\mathfrak b$, where $\mathfrak q$ is a
prime ideal of $C$, and $C'_{\mathfrak q'}$ is isomorphic to
$C_{\mathfrak q}/\mathfrak bC_{\mathfrak q}$.
\end{proof}

\subsection{Morphisms locally of finite presentation}
\label{subsection:IV.1.4}

\begin{env}[1.4.1]
\label{IV.1.4.1}
Given a ring $A$, we will say that a (commutative) $A$-algebra $B$
is \emph{of finite presentation} if it is isomorphic to the quotient
$A[T_1,\ldots,T_n]/\mathfrak a$ of a polynomial algebra over $A$ by an
ideal $\mathfrak a$ \emph{of finite type} of $A[T_1,\ldots,T_n]$. For every
(commutative) $A$-algebra $A'$, $B\otimes_AA'$ is then an $A'$-algebra of
finite presentation, being isomorphic to the quotient of $A'[T_1,\ldots,T_n]$
by the canonical image in this ring of $\mathfrak a\otimes_AA'$, which is
evidently an $A'[T_1,\ldots,T_n]$-module of finite type. If $B$ is an
$A$-algebra of finite presentation and $C$ a $B$-algebra of finite
presentation, then $C$ is an $A$-algebra of finite presentation. Indeed, let
$B=A[S_1,\ldots,S_m]/\mathfrak a$ and
$C=B[T_1,\ldots,T_n]/\mathfrak b$, where $\mathfrak a$ (resp. $\mathfrak b$) is
an ideal of finite type of $A[S_1,\ldots,S_m]$ (resp.
$B[T_1,\ldots,T_n]$); the ring $B[T_1,\ldots,T_n]$ is isomorphic to
$A[S_1,\ldots,S_m,T_1,\ldots,T_n]/\mathfrak a'$, where $\mathfrak a'$ is
the canonical image of
$\mathfrak a\otimes_AA[T_1,\ldots,T_n]$, and is therefore an ideal of finite
type of $A[S_1,\ldots,S_m,T_1,\ldots,T_n]$; the ideal $\mathfrak b$ of
$B[T_1,\ldots,T_n]$ is of the form $\mathfrak b'/\mathfrak a'$, where
$\mathfrak b'$ is an ideal of $A[S_1,\ldots,S_m,T_1,\ldots,T_n]$; since
$\mathfrak a'$ and $\mathfrak b'/\mathfrak a'$ are modules of finite type
over $A[S_1,\ldots,S_m,T_1,\ldots,T_n]$, so is $\mathfrak b'$,
and $C$, being isomorphic to
$A[S_1,\ldots,S_m,T_1,\ldots,T_n]/\mathfrak b'$, is therefore indeed an
$A$-algebra of finite presentation. Finally, note that if $A$ is
\emph{Noetherian}, so is $A[T_1,\ldots,T_n]$, hence every
$A$-algebra of finite type is then of finite presentation.
\end{env}

\begin{definition}[1.4.2]
\label{IV.1.4.2}
Let $f:X\to Y$ be a morphism of preschemes, $x$ a point of $X$, $y=f(x)$.
We say that $f$ is \emph{of finite presentation at $x$} if there exist an
affine open neighbourhood $V$ of $y$ and an affine open neighbourhood $U$ of
$x$ such that $f(U)\subset V$ and the ring $A(U)$ is an $A(V)$-algebra of
finite presentation. We say that $f$ is \emph{locally of finite presentation}
if it is of finite presentation at every point of $X$.
\end{definition}

If $Y$ is locally Noetherian, it amounts to the same to say that $f$ is
locally of finite type or locally of finite presentation.

\begin{proposition}[1.4.3]
\label{IV.1.4.3}
\begin{enumerate}
  \item[(i)] Every local isomorphism is locally of finite presentation.
  \item[(ii)] If two morphisms $f:X\to Y$, $g:Y\to Z$ are locally of
  finite presentation, then so is $g\circ f$.

\oldpage[IV-1]{231}

  \item[(iii)] If $f:X\to Y$ is an $S$-morphism locally of finite
  presentation, $f_{(S')}:X_{(S')}\to Y_{(S')}$ is locally of finite
  presentation for every extension $S'\to S$ of the base prescheme.
  \item[(iv)] If $f:X\to X'$ and $g:Y\to Y'$ are two $S$-morphisms
  locally of finite presentation, $f\times_Sg$ is locally of finite
  presentation.
  \item[(v)] If the composite $g\circ f$ of two morphisms is locally of
  finite presentation and if $g$ is locally of finite type, then $f$ is
  locally of finite presentation.
\end{enumerate}
\end{proposition}

\begin{proof}
Assertion (i) is trivial. To prove (iii), we may restrict to the case
where $S=Y$ (I, 3.3.11); let $x'$ be a point of $X_{(S')}$, $s'$ and $x$ its
projections onto $S'$ and $X$ respectively, $s$ the common projection of $s'$
and $x$ onto $S$. By hypothesis, there is an affine open neighbourhood $V$
(resp. $U$) of $s$ (resp. $x$) such that the image of $U$ is contained in $V$
and $A(U)$ is an $A(V)$-algebra of finite presentation; let $V'$ be an affine
open neighbourhood of $s'$ whose image in $S$ is contained in $V$; then
$U\times_VV'$ is an affine open neighbourhood of $x'$ whose image in $S'$
is contained in $V'$ and whose ring $A(U)\otimes_{A(V)}A(V')$ is an
$A(V')$-algebra of finite presentation (1.4.1).

To prove (ii), consider a point $x\in X$, and let $y=f(x)$,
$z=g(f(x))$. By hypothesis there is an affine open neighbourhood $W=\operatorname{Spec}(C)$
(resp. $V=\operatorname{Spec}(B)$) of $z$ (resp. $y$) such that $g(V)\subset W$ and $B$
is a $C$-algebra of finite presentation. On the other hand, there is an affine
open neighbourhood $V'=\operatorname{Spec}(B')$ (resp. $U=\operatorname{Spec}(A)$) of $y$ (resp. $x$) such that
$f(U)\subset V'$ and $A$ is a $B'$-algebra of finite presentation. Let
$s\in B$ be such that the open subset $D(s)=\operatorname{Spec}(B_s)$ of $V$ is a neighbourhood of $y$
contained in $V'$, and let $U'=f^{-1}(D(s))\subset U$; we have
$U'=\operatorname{Spec}(A\otimes_{B'}B_s)$, and $A\otimes_{B'}B_s$ is a $B_s$-algebra of
finite presentation (1.4.1); on the other hand $B_s=B[T]/(1-sT)$ is a
$B$-algebra of finite presentation by definition; consequently (1.4.1),
$A\otimes_{B'}B_s$ is a $C$-algebra of finite presentation, which proves
(ii).

We know that (iv) follows from (i), (ii) and (iii) (I, 3.5.1). To prove (v),
consider the commutative diagram
\[
\xymatrix{
X \ar[r]^{\Gamma_f} \ar[d] & X\times_ZY \ar[d]^{f\times 1} \\
Y \ar[r]_{\Delta_g} & Y\times_ZY
}
\]
which identifies $X$ with the product of the $(Y\times_ZY)$-preschemes $Y$ and
$X\times_ZY$ (I, 5.3.7, $\Delta_g$ being the diagonal morphism). If we know
that $\Delta_g$ is locally of finite presentation, it will follow by
(iii) that so is $\Gamma_f$. On the other hand, we have the factorization
of $f$ as
$X\xrightarrow{\Gamma_f}X\times_ZY\xrightarrow{p_2}Y$ (I, 5.3.13), where the
projection $p_2$ is equal to $(g\circ f)\times_Z1_Y$; since $g\circ f$ is
assumed to be locally of finite presentation, so is $p_2$ by
(iv), and it will finally follow from (ii) that $f$ is also locally of
finite presentation. We therefore see that we are reduced to proving the
\end{proof}

\begin{corollary}[1.4.3.1]
\label{IV.1.4.3.1}
Let $g:Y\to Z$ be a morphism locally of finite type; then the diagonal
morphism $\Delta_g:Y\to Y\times_ZY$ is locally of finite presentation.
\end{corollary}

\oldpage[IV-1]{232}

\begin{proof}
We may restrict to the case where $Z=\operatorname{Spec}(A)$, $Y=\operatorname{Spec}(B)$, $B$ being an
$A$-algebra of finite type; the morphism $\Delta_g$ then corresponds to
the \emph{augmentation homomorphism} $\delta:B\otimes_AB\to B$ such that
$\delta(x\otimes y)=xy$. Taking into account Definition (1.4.2), it suffices
to show that the kernel $\mathfrak J$ of $\delta$ is an ideal \emph{of finite
type}, which follows from (0, 20.4.4).
\end{proof}

\begin{proposition}[1.4.4]
\label{IV.1.4.4}
Let $A$ be a ring, $B$ an $A$-algebra, $B'$ a polynomial algebra
$A[T_1,\ldots,T_n]$, $u:B'\to B$ a surjective homomorphism of
$A$-algebras. For $B$ to be an $A$-algebra of finite presentation, it is
necessary and sufficient that the kernel $\mathfrak J$ of $u$ be an ideal of
finite type of $B'$.
\end{proposition}

\begin{proof}
The condition is sufficient by definition (1.4.1). To see that it is
necessary, observe that the morphism $g:\operatorname{Spec}(B')\to\operatorname{Spec}(A)$ is
locally of finite type; if $B$ is an $A$-algebra of finite presentation,
the morphism $f:\operatorname{Spec}(B)\to\operatorname{Spec}(B')$ corresponding to $u$ is such that
$g\circ f:\operatorname{Spec}(B)\to\operatorname{Spec}(A)$ is locally of finite presentation, hence it
follows from (1.4.3, (v)) that $f$ is also locally of finite presentation.
Now, to prove that the ideal $\mathfrak J$ is of finite type, it suffices to
prove that $\widetilde{\mathfrak J}$ is a $\widetilde{B'}$-module of finite
type (II, 6.1.4.1). Going back to Definition (1.4.2), we are finally reduced
to proving the
\end{proof}

\begin{lemma}[1.4.4.1]
\label{IV.1.4.4.1}
Let $\mathfrak a$ be an ideal of a ring $C$, $s$ an element of $C$, such that
$C_s/\mathfrak a_s$ is a $C$-algebra of finite presentation; then
$\mathfrak a_s$ is an ideal of finite type in the ring $C_s$.
\end{lemma}

\begin{proof}
By hypothesis, there exist a polynomial $C$-algebra
$C'=C[T_1,\ldots,T_n]$ and a surjective $C$-homomorphism
$v:C'\to C_s/\mathfrak a_s$ whose kernel $\mathfrak b$ is of finite type.
Let $p:C_s\to C_s/\mathfrak a_s$ be the canonical homomorphism; for every $i$,
there is a $t_i\in C$ such that $p(t_i/s^k)=v(T_i)$ (we may suppose that
the exponent $k$ is the same for all the indices $i$). Then consider the
$C_s$-algebra $D=C_s[T_1,\ldots,T_n]$, and the homomorphism of $C_s$-algebras
$w:D\to C_s$ such that $w(T_i)=t_i/s^k$; $w$ is evidently surjective, and
consequently so is $v'=p\circ w:D\to C_s/\mathfrak a_s$. On the other
hand, every polynomial of $D$ can be written $P/s^m$, where
$P\in C'=C[T_1,\ldots,T_n]$, and we have
$v'(P/s^m)=(1/s^m)v(P)$; since $1/s$ is invertible in $C_s$, the relation
$v'(P/s^m)=0$ in $C_s/\mathfrak a_s$ is equivalent to $v(P)=0$, and
consequently the kernel $\mathfrak b'$ of $v'$ is generated by the canonical
image of $\mathfrak b$ in $D$, and \emph{a fortiori} is of finite type in the
ring $D$. But
$\mathfrak b'=v'^{-1}(0)=w^{-1}(p^{-1}(0))=w^{-1}(\mathfrak a_s)$, and since
$w$ is surjective, $\mathfrak a_s=w(\mathfrak b')$, which proves that
$\mathfrak a_s$ is an ideal of finite type of $C_s$.
\end{proof}

\begin{corollary}[1.4.5]
\label{IV.1.4.5}
Let $X$, $Y$ be two preschemes, $j:X\to Y$ an immersion, $U$ an open subset
of $Y$ such that $j(X)$ is closed in $U$, $\mathcal J$ the quasi-coherent
ideal of $\mathcal O_U$ that defines the closed sub-prescheme of $Y$
associated to $j$. For $j$ to be locally of finite presentation, it is
necessary and sufficient that $\mathcal J$ be an $\mathcal O_U$-module of
finite type.
\end{corollary}

\par\medskip\noindent
\begin{proposition}[1.4.6]
\label{IV.1.4.6}
Let $\varphi:A\to B$ be a ring homomorphism. For the corresponding
morphism $f:\operatorname{Spec}(B)\to\operatorname{Spec}(A)$ to be locally of finite presentation,
it is necessary and sufficient that $B$ be an $A$-algebra of finite
presentation.
\end{proposition}

\begin{proof}
The condition being trivially sufficient, it remains to prove that it is
necessary. If $f$ is locally of finite presentation, it already follows from
(1.3.6) that $B$ is an $A$-algebra of finite type, hence there exists a
surjective homomorphism $u:B'=A[T_1,\ldots,T_n]\to B$ of $A$-algebras; it
follows from (1.4.3, (v)) that the immersion
$j:\operatorname{Spec}(B)\to\operatorname{Spec}(B')$ is locally

\oldpage[IV-1]{233}

of finite presentation; hence (1.4.5), if $\mathfrak J$ is the kernel of $u$,
the $\widetilde{B'}$-module $\widetilde{\mathfrak J}$ is of finite type, and
consequently (II, 6.1.4.1) $\mathfrak J$ is an ideal of finite type.
\end{proof}

\begin{proposition}[1.4.7]
\label{IV.1.4.7}
Let $\rho:A\to B$ be a ring homomorphism making $B$ a finite
$A$-algebra. For $B$ to be an $A$-algebra of finite presentation, it is
necessary and sufficient that $B$ be an $A$-module of finite presentation.
\end{proposition}

We will use the

\begin{lemma}[1.4.7.1]
\label{IV.1.4.7.1}
If $B$ is a finite $A$-algebra, there exists a surjective $A$-homomorphism
of algebras $u:B'\to B$, where $B'$ is a finite $A$-algebra of finite
presentation which is a free $A$-module.
\end{lemma}

\begin{proof}
Indeed, there is a finite system $(a_i)_{1\leq i\leq m}$ of generators of the
$A$-algebra $B$ such that each $a_i$ satisfies a relation $F_i(a_i)=0$, where
$F_i\in A[X]$ is a monic polynomial of degree $>0$; the quotient $A$-algebra
$B'_i=A[X]/(F_i)$ is therefore free of finite rank over $A$; let $c_i$ be the
class of $X$ in $B'_i$. There exists an $A$-homomorphism of algebras
$u_i:B'_i\to B$ such that $u_i(c_i)=a_i$; it then suffices to take for $B'$
the tensor product $B'_1\otimes_AB'_2\otimes_A\cdots\otimes_AB'_m$ and to
consider the homomorphism $u_1\otimes u_2\otimes\cdots\otimes u_m:B'\to B$,
which is surjective by construction. Moreover, if $B''=A[T_1,\ldots,T_m]$
and if $b'_i$ denotes the canonical image of $c_i$ in $B'$, the
$A$-homomorphism of algebras $v:B''\to B'$ such that $v(T_i)=b'_i$
($1\leq i\leq m$) is surjective, and its kernel $\mathfrak b'$ is the ideal
of $B''$ generated by the $F_i(T_i)$, hence of finite type; this proves that
$B'$ is an $A$-algebra of finite presentation.
\end{proof}

\begin{proof}
This lemma having been proved, let us keep the same notation, and let
$w=v\circ u:B''\to B$. If $\mathfrak b$ (resp. $\mathfrak a$) is the kernel of
$w$ (resp. $u$), we have $\mathfrak a=v(\mathfrak b)$ since $v$ is
surjective. Now (1.4.4), if $B$ is an $A$-algebra of finite presentation,
$\mathfrak b$ is an ideal of finite type of $B''$, hence $\mathfrak a$ is an
ideal of finite type of $B'$, hence an $A$-module of finite type since $B'$ is
a finite $A$-algebra; since $B'$ is a free $A$-module, $B$ is an $A$-module of
finite presentation. Conversely, if $B'$ is an $A$-module of finite
presentation, $\mathfrak a$ is an $A$-module of finite type (Bourbaki,
\emph{Alg. comm.}, chap. I, \S~2, n\textsuperscript{o}~8, lemma 9) and
\emph{a fortiori} an ideal of finite type of $B'$; consequently, $B$ is by
definition a $B'$-algebra of finite presentation, and since $B'$ is an
$A$-algebra of finite presentation, $B$ is an $A$-algebra of finite
presentation.
\end{proof}

\subsection{Morphisms of finite type}
\label{subsection:IV.1.5}

\begin{proposition}[1.5.1]
\label{IV.1.5.1}
Let $f:X\to Y$ be a morphism, $(U_\alpha)$ a covering of $Y$ consisting of
affine open subsets. The following conditions are equivalent:
\begin{enumerate}
  \item[a)] $f$ is locally of finite type and quasi-compact.
  \item[b)] For every $\alpha$, $f^{-1}(U_\alpha)$ is a finite union of
  affine open subsets $V_{\alpha i}$ such that each ring $A(V_{\alpha i})$ is
  an $A(U_\alpha)$-algebra of finite type.
  \item[c)] For every affine open subset $U$ of $Y$, $f^{-1}(U)$ is a finite
  union of affine open subsets $V_j$ such that the rings $A(V_j)$ are
  $A(U)$-algebras of finite type.
\end{enumerate}
\end{proposition}

For the proof, see (I, 6.3.2 and 6.6.3).

\begin{definition}[1.5.2]
\label{IV.1.5.2}
We say that a morphism $f$ is \emph{of finite type} if it satisfies the
equivalent conditions of Proposition (1.5.1).
\end{definition}

\oldpage[IV-1]{234}

\begin{proposition}[1.5.3]
\label{IV.1.5.3}
Let $f:X\to Y$ be a morphism of finite type; if $Y$ is Noetherian, so is
$X$.
\end{proposition}

For the proof, see (I, 6.3.7).

\begin{proposition}[1.5.4]
\label{IV.1.5.4}
\begin{enumerate}
  \item[(i)] Let $j:X\to Y$ be an immersion morphism. If $j$ is
  quasi-compact (which is the case if $j$ is a closed immersion, or if
  the underlying space of $X$ is Noetherian, or if the underlying space of
  $Y$ is locally Noetherian), $j$ is of finite type.
  \item[(ii)] The composite of two morphisms of finite type is of finite type.
  \item[(iii)] If $f:X\to Y$ is an $S$-morphism of finite type,
  $f_{(S')}:X_{(S')}\to Y_{(S')}$ is of finite type for every extension
  $S'\to S$ of the base prescheme.
  \item[(iv)] If $f:X\to X'$ and $g:Y\to Y'$ are two $S$-morphisms of finite
  type, $f\times_Sg$ is of finite type.
  \item[(v)] If the composite $g\circ f$ of two morphisms is of finite type
  and if $g$ is quasi-separated or if $X$ is Noetherian, $f$ is of finite type.
  \item[(vi)] If a morphism $f$ is of finite type, then so is
  $f_{\mathrm{red}}$.
\end{enumerate}
\end{proposition}

This follows immediately (except for case (v) where $X$ is Noetherian, proved
in (I, 6.3.6)) from the definitions and from (1.1.2), (1.2.4) and (1.3.4).

\begin{corollary}[1.5.5]
\label{IV.1.5.5}
Let $f:X\to Y$ be a morphism of finite type. For every morphism $Y'\to Y$
such that $Y'$ is Noetherian, $X\times_YY'$ is Noetherian.
\end{corollary}

This follows from (1.5.4, (iii)) and from (1.5.3).

\begin{corollary}[1.5.6]
\label{IV.1.5.6}
Let $X$ be a prescheme of finite type over a locally Noetherian prescheme $S$.
Then every $S$-morphism $f:X\to Y$ is of finite type.
\end{corollary}

For the proof, see (I, 6.3.9).

\begin{proposition}[1.5.7]
\label{IV.1.5.7}
Let $\varphi:A\to B$ be a ring homomorphism. For the corresponding
morphism $f:\operatorname{Spec}(B)\to\operatorname{Spec}(A)$ to be of finite type, it is necessary and
sufficient that $\varphi$ make $B$ an $A$-algebra of finite type.
\end{proposition}

For the proof, see (I, 6.3.3).

\subsection{Morphisms of finite presentation}
\label{subsection:IV.1.6}

\begin{definition}[1.6.1]
\label{IV.1.6.1}
Let $X$, $Y$ be two preschemes. We say that a morphism $f:X\to Y$ is
\emph{of finite presentation} if it satisfies the following three conditions:
\begin{enumerate}
  \item[(i)] $f$ is locally of finite presentation;
  \item[(ii)] $f$ is quasi-compact (which, when (i) holds, is equivalent
  to saying that $f$ is of finite type (1.5.2));
  \item[(iii)] $f$ is quasi-separated.
\end{enumerate}
We then also say that $X$ is \emph{of finite presentation over $Y$}, or a
\emph{$Y$-prescheme of finite presentation}.
\end{definition}

Condition (iii) is automatically satisfied if $f$ is separated, or if $X$
is locally Noetherian (1.2.8); when $Y$ is locally Noetherian, it
amounts to the same to say that $f$ is \emph{of finite presentation} or that
$f$ is \emph{of finite type} (the latter condition implying that $X$ is
locally Noetherian (1.3.3)).

\oldpage[IV-1]{235}

\begin{proposition}[1.6.2]
\label{IV.1.6.2}
\begin{enumerate}
  \item[(i)] Every quasi-compact immersion that is locally of finite
  presentation (in particular every quasi-compact open immersion) is of
  finite presentation.
  \item[(ii)] The composite of two morphisms of finite presentation is of
  finite presentation.
  \item[(iii)] Let $X$, $Y$ be two $S$-preschemes, $f:X\to Y$ an
  $S$-morphism of finite presentation. Then, for every morphism $S'\to S$,
  the morphism $f_{(S')}:X_{(S')}\to Y_{(S')}$ is of finite presentation.
  \item[(iv)] Let $f:X\to Y$, $f':X'\to Y'$ be two $S$-morphisms of finite
  presentation; then
  $f\times_Sf':X\times_SX'\to Y\times_SY'$ is of finite presentation.
  \item[(v)] If the composite $g\circ f$ of two morphisms is of finite
  presentation and if $g$ is quasi-separated and locally of finite
  presentation, $f$ is of finite presentation.
\end{enumerate}
\end{proposition}

This follows immediately from (1.1.2), (1.2.2), (1.2.4) and (1.4.3).

It follows in particular from (1.6.2, (iii)) that if $f$ is a morphism of
finite presentation and $U$ an open subset of $Y$, the restriction
$f^{-1}(U)\to U$ of $f$ is a morphism of finite presentation. Conversely,
let $(U_\alpha)$ be a covering of $Y$ by affine open subsets, and suppose
that the restrictions $f^{-1}(U_\alpha)\to U_\alpha$ of $f$ are
morphisms of finite presentation; it then follows that $f$ is of finite
presentation, since $f$ is evidently locally of finite presentation,
quasi-compact by virtue of (1.1.1) and quasi-separated by virtue of (1.2.6).
We can therefore say that the property, for a morphism $f:X\to Y$, of being
of finite presentation is \emph{local on $Y$}.

If $X$ is a quasi-separated prescheme, every morphism $f:X\to Y$ is
quasi-separated (1.2.3). Hence, if $f$ is quasi-compact and locally of finite
presentation, $f$ is then of finite presentation.

In particular:

\begin{corollary}[1.6.3]
\label{IV.1.6.3}
Let $\varphi:A\to B$ be a ring homomorphism. For the corresponding
morphism $f:\operatorname{Spec}(B)\to\operatorname{Spec}(A)$ to be of finite presentation, it is necessary and
sufficient that $\varphi$ make $B$ an $A$-algebra of finite presentation.
\end{corollary}

Indeed, the necessity of the condition follows from (1.4.6).

\begin{remark}[1.6.4]
\label{IV.1.6.4}
In Definition (1.6.1), condition (iii) is not a consequence of the other
two. For example, let $Y$ be an affine scheme whose underlying space is not
Noetherian, and let $U$ be an open subset of $Y$ that is not quasi-compact.
Let $X$ be the prescheme obtained by gluing two preschemes $Y_1$, $Y_2$
isomorphic to $Y$ along the open subsets $U_1$, $U_2$ corresponding to $U$
(0$_{\mathrm I}$,
4.1.7), so that $X$ is the union of two open subsets isomorphic respectively
to $Y_1$ and $Y_2$ (and which we identify with the latter), with
$Y_1\cap Y_2=U$. Moreover, there is a morphism $f:X\to Y$ that coincides on
$Y_i$ with the given isomorphism $Y_i\to Y$ ($i=1,2$). It is clear that $f$
satisfies condition (i) of (1.6.1), being a local isomorphism; it also
satisfies (ii), the inverse image of a quasi-compact open subset of $Y$ being
the union of its images in $Y_1$ and $Y_2$; but since $Y_1\cap Y_2$ is not
quasi-compact, $f$ is not quasi-separated (1.2.7).
\end{remark}

\begin{proposition}[1.6.5]
\label{IV.1.6.5}
Let $X'$, $X''$ be two preschemes, $X=X'\amalg X''$ their sum,
$f':X'\to Y$, $f'':X''\to Y$ two morphisms, $f:X\to Y$ the morphism that
coincides with $f'$ on $X'$ and with $f''$ on $X''$. For $f$ to be of
finite presentation, it is necessary and sufficient that $f'$ and $f''$ be so.
\end{proposition}

\par\medskip\noindent
\begin{proof}
It suffices to show that for $f$ to have one of the three properties of
Definition

\oldpage[IV-1]{236}

(1.6.1), it is necessary and sufficient that $f'$ and $f''$ have it; this is
evident for the property of being locally of finite presentation, which is
local on $X$; for the property of being quasi-compact, this was seen in
(1.1.6), and for the property of being quasi-separated, in (1.2.9).
\end{proof}


\subsection{Improvements of earlier results}
\label{subsection:IV.1.7}

In this subsection we give a list of propositions proved in the preceding
chapters whose statements can be improved with the help of the new
finiteness conditions introduced above.

\begin{env}[1.7.1]
\label{IV.1.7.1}
In the statements of (I, 6.4.2, 6.4.3 and 6.4.9), instead of supposing
that $X$ and $Y$ are of finite type over the field $K$, we may suppose only
that they are \emph{locally of finite type} over $K$; for (I, 6.4.2), it
suffices to observe that every point of a prescheme $X$ that is closed in
every affine open subset of $X$ is closed in $X$, and an affine scheme locally
of finite type over $K$ is \emph{ipso facto} of finite type (1.5.1). For
(I, 6.4.9), we use (1.3.7) and (1.3.4, (v)).
\end{env}

\begin{env}[1.7.2]
\label{IV.1.7.2}
In (I, 6.5.1, (ii)), instead of supposing that $S$ is locally
Noetherian and $Y$ of finite type over $S$, we may suppose only that $Y$ is
\emph{locally of finite presentation} over $S$, as the proof immediately
shows (by applying Definition (1.4.2)). Likewise in (I, 6.5.4,
(ii)) and (I, 6.5.5, (ii)) it suffices to suppose that $f$ is an
$S$-morphism, $Y$ being locally of finite presentation over $S$.
\end{env}

\begin{env}[1.7.3]
\label{IV.1.7.3}
In (I, 7.1.11), for the second assertion, instead of supposing $S$
locally Noetherian and $Y$ of finite type over $S$, we may suppose only
$Y$ locally of finite presentation over $S$ (the proof depending on
(I, 6.5.1, (ii))); likewise in (I, 7.1.12 to 7.1.15).
\end{env}

\begin{env}[1.7.4]\label{IV.1.7.4}
\EGAReaderTarget{ega4.c04.statement.1-7-4}{IV.1.7.4-fr}
In (I, 9.2.2), we may replace the three hypotheses by the single hypothesis
(implied by each of the others) that $f$ is \emph{quasi-compact and
quasi-separated}, by virtue of (1.2.6) and (1.2.7). Note that in
(I, 9.2.1), the hypothesis means exactly that $f$ is quasi-compact and
quasi-separated.
\end{env}

\begin{env}[1.7.5]
\label{IV.1.7.5}
In (I, 9.3.1, 9.3.2 and 9.3.3), we may replace the two hypotheses on $X$
by the (weaker) hypothesis that $X$ is a \emph{quasi-compact and
quasi-separated} prescheme, the proof using exactly these two properties (by
virtue of (1.2.7)).
\end{env}

\begin{env}[1.7.6]
\label{IV.1.7.6}
In (I, 9.3.4 and 9.3.5), it suffices to suppose $X$ quasi-compact,
$\mathcal F$ and $\mathcal J$ quasi-coherent and of finite type.
\end{env}

\begin{env}[1.7.7]
\label{IV.1.7.7}
In (I, 9.4.7), we may weaken hypothesis b) by supposing only $X$
\emph{quasi-separated}, since it suffices only that the canonical immersion
$U\to X$ be quasi-compact (1.2.7). We may make the same modification in
(I, 9.4.8, 9.4.9 and 9.4.10).
\end{env}

\begin{env}[1.7.8]
\label{IV.1.7.8}
In (I, 9.5.1), the conditions indicated in parentheses in the statement
may be replaced by the condition that $f$ be quasi-compact and
quasi-separated.
\end{env}

\begin{env}[1.7.9]
\label{IV.1.7.9}
In (I, 9.6.5), we may replace the hypotheses on $X$ by the weaker
hypothesis that $X$ is quasi-compact and quasi-separated, as the
proof shows,

\oldpage[IV-1]{237}

which uses only (I, 9.4.7). In (I, 9.6.6), the hypothesis ``$X$ is a
quasi-compact scheme'' may be replaced by ``$X$ is a quasi-compact and
quasi-separated prescheme''.
\end{env}

\begin{env}[1.7.10]
\label{IV.1.7.10}
In (II, 1.7.15), we may suppose only that $Y$ is quasi-compact and
quasi-separated, the proof using only (I, 9.6.5).
\end{env}

\begin{env}[1.7.11]
\label{IV.1.7.11}
In the second assertion of (II, 3.4.5), we may substitute for the two
hypotheses on $Y$ the weaker hypothesis that $Y$ is quasi-compact and
quasi-separated. Likewise in (II, 3.4.8).
\end{env}

\begin{env}[1.7.12]
\label{IV.1.7.12}
In (II, 3.8.5), we may likewise suppose only that $Y$ is quasi-compact
and quasi-separated, this hypothesis entering through (I, 9.4.7).
\end{env}

\begin{env}[1.7.13]
\label{IV.1.7.13}
In (II, 4.4.1, 4.4.6 and 4.4.7), it suffices to suppose that $Y$ is
\emph{quasi-compact and quasi-separated}, taking into account (1.2.3) in the
proof of (II, 4.4.7). Likewise, in (II, 4.4.10.1), it suffices to
suppose $Z$ quasi-compact and quasi-separated.
\end{env}

\begin{env}[1.7.14]
\label{IV.1.7.14}
In (II, 4.5.2) and (II, 4.5.5), it suffices to suppose $X$
\emph{quasi-compact and quasi-separated}, this hypothesis entering through (I,
9.3.1 and 9.3.2).
\end{env}

\begin{env}[1.7.15]
\label{IV.1.7.15}
In (II, 4.6.8), it again suffices to suppose $X$ \emph{quasi-compact and
quasi-separated}, this hypothesis entering through (I, 9.4.7).
\end{env}

\begin{env}[1.7.16]
\label{IV.1.7.16}
In (II, 5.1.2), it suffices to suppose that $X$ is quasi-compact and
quasi-separated (the hypothesis entering through (II, 4.5.2 and 4.5.5)).
Likewise in (II, 5.1.9), where we use (I, 9.6.5).
\end{env}

\begin{env}[1.7.17]
\label{IV.1.7.17}
In (II, 5.2.1), it still suffices to suppose $X$ quasi-compact and
quasi-separated (the hypothesis entering through (II, 4.5.2)).
\end{env}

\begin{env}[1.7.18]
\label{IV.1.7.18}
In (II, 5.2.2), one must suppose $f$ quasi-compact and $X$ quasi-separated;
the present proof is indeed insufficient (with the hypotheses made),
because from the fact that for every quasi-coherent $\mathcal O_X$-module
$\mathcal F$ we have $\mathrm R^1f_*(\mathcal F)=0$, it does not necessarily
follow, for an affine open subset $U$ of $Y$, that the same property holds
for the quasi-coherent $(\mathcal O_X|f^{-1}(U))$-modules, if we do not know
that such a module is the \emph{restriction} of a quasi-coherent
$\mathcal O_X$-module (the same remark applying to conditions b) and c'));
when $f$ is quasi-compact and $X$ quasi-separated, this extension is possible
by virtue of (I, 9.4.7), modified above (1.7.7).
\end{env}

\begin{env}[1.7.19]
\label{IV.1.7.19}
In (II, 5.3.2), it suffices to suppose $Y$ \emph{quasi-compact and
quasi-separated} (which enters through (II, 4.4.6 and 4.4.7)). Likewise in
(II, 5.5.3, (ii)).
\end{env}

\begin{env}[1.7.20]
\label{IV.1.7.20}
In (III, 1.2.2, 1.2.3 and 1.2.4), we may replace the two hypotheses on
$X$ by the weaker hypothesis that $X$ is a \emph{quasi-compact
and quasi-separated} prescheme, the proof using only (I, 9.3.3).
\end{env}

\begin{env}[1.7.21]
\label{IV.1.7.21}
In (III, 1.4.10 to 1.4.14), we may replace ``separated'' by
``quasi-separated'', this hypothesis serving only to make it possible to
apply (I, 9.2.2) (see (1.7.4) above). Likewise, in (III, 1.4.15), it
suffices to suppose that the morphism $u$ is \emph{quasi-separated and
quasi-compact}.
\end{env}

\begin{env}[1.7.22]
\label{IV.1.7.22}
In (III, 2.1.3), we may again replace the hypotheses by the weaker
hypothesis that $X$ is a quasi-compact and quasi-separated prescheme, the
argument being the same as in (III, 1.2.2).
\end{env}

\subsection{Morphisms of finite presentation and constructible sets}
\label{subsection:IV.1.8}

\oldpage[IV-1]{238}

\begin{env}[1.8.1]
\label{IV.1.8.1}
Let $X$ be a prescheme; we know that every quasi-compact open subset of $X$
is a finite union of affine open subsets, and conversely. For an open subset
$U$ of $X$ to be \emph{retrocompact} in $X$ (0\textsubscript{III}, 9.1.1), it
is therefore necessary and sufficient that for every \emph{affine open subset}
$V$ of $X$, $U\cap V$ be quasi-compact. When $X$ is \emph{quasi-separated}
(1.2.1), every quasi-compact open subset of $X$ is retrocompact in $X$
(1.2.7), hence every locally constructible subset of $X$ is retrocompact in
$X$ (0\textsubscript{III}, 9.1.13); if \emph{moreover} $X$ is
\emph{quasi-compact}, the constructible subsets and the locally constructible
subsets of $X$ are the same (0\textsubscript{III}, 9.1.12), and so are the
constructible open subsets and the quasi-compact open subsets
(0\textsubscript{III}, 9.1.5).
\end{env}

\begin{proposition}[1.8.2]
\label{IV.1.8.2}
Let $X$, $Y$ be two preschemes, $f:X\to Y$ a morphism. For every
constructible (resp. locally constructible) subset $Z$ of $Y$, $f^{-1}(Z)$ is
a constructible (resp. locally constructible) subset of $X$.
\end{proposition}

\begin{proof}
Suppose $Z$ locally constructible; for every $x\in X$, there will be an
affine open neighbourhood $V$ of $f(x)$ such that $Z\cap V$ is constructible
in $V$; if the proposition is proved for constructible subsets,
$f^{-1}(Z)\cap f^{-1}(V)$ will be constructible in $f^{-1}(V)$, hence
$f^{-1}(Z)$ will be locally constructible. It therefore suffices to consider
the case where $Z$ is constructible, and taking into account
(0\textsubscript{III}, 9.1.3), we are reduced to the case where $Z$ is open
and \emph{retrocompact} in $Y$, in other words such that the canonical
injection $Z\to Y$ is a quasi-compact morphism; it then suffices to see that
$f^{-1}(Z)$ is \emph{retrocompact in $X$}, in other words such that the
canonical injection $f^{-1}(Z)\to X$ is a quasi-compact morphism; but this
follows from (1.1.2, (iii)), since $f^{-1}(Z)=Z\times_YX$ (I, 4.4.1).
\end{proof}

\begin{lemma}[1.8.3]
\label{IV.1.8.3}
Let $X$ be a quasi-compact and quasi-separated prescheme, $Z$ a
constructible subset of $X$. Then there exist an affine scheme $X'$ and a
morphism of finite presentation $f:X'\to X$ such that $f(X')=Z$.
\end{lemma}

\begin{proof}
Since $X$ is quasi-compact, it is a finite union of affine open subsets $X_i$
($i\in I$), and since $X$ is quasi-separated, it follows from (1.2.7) that
each of the canonical immersions $g_i:X_i\to X$ is of finite presentation; we
conclude that if $f_i:X'_i\to X_i$ is a morphism of finite presentation, so
is $g_i\circ f_i:X'_i\to X$ (1.6.2), and if $X'$ is the sum prescheme of the
$X'_i$, the morphism $h:X'\to X$ that coincides with $g_i\circ f_i$ on each
$X'_i$ is also of finite presentation (1.6.5). Since $Z\cap X_i$ is
constructible in $X_i$ (0\textsubscript{III}, 9.1.8), we see that we are
reduced to proving the lemma when $X_i$ is \emph{affine} with ring $A$. Since
$Z$ is then a finite union of sets of the form $U\cap\complement V$, where
$U$ and $V$ are quasi-compact open subsets (0\textsubscript{III}, 9.1.3), we
see, by again considering a suitable sum of preschemes, that we may reduce to
the case where $Z=U\cap\complement V$; since moreover $U$ is a finite union
of affine open subsets, we may even restrict to the case where $U$ is
\emph{affine}. Since $V$ is quasi-compact, it is a finite union of open
subsets of the form $X_{f_i}$, where $f_i\in A$; let $\mathfrak J$ be the
ideal of $A$ generated by the $f_i$, and let $Z'$ be the affine closed
subscheme of $X$ that it defines (and which is by construction of finite
presentation over $X$); we have by definition $V=\complement Z'$, hence
$U\cap\complement V=U\cap Z'$. Consider the affine scheme $X'=Z'\times_XU$

\oldpage[IV-1]{239}

and let $f:X'\to X$ be the structural morphism; we have just seen that
the canonical immersion $Z'\to X$ is of finite presentation, and so is
the open immersion $U\to X$, which is quasi-compact since $U$ and
$X$ are affine; we therefore conclude from (1.6.2, (iv)) that $f$ is of
finite presentation, and we have $f(X')=Z'\cap U$ (I, 3.4.8), which completes
the proof.
\end{proof}

\begin{theorem}[1.8.4]
\label{IV.1.8.4}
\textup{(Chevalley).} --- Let $f:X\to Y$ be a morphism of finite
presentation. For every locally constructible subset $Z$ of $X$, $f(Z)$ is
locally constructible in $Y$.
\end{theorem}

\begin{proof}
Let $y\in Y$ and let $V$ be an affine open neighbourhood of $y$; since $f$ is
quasi-compact and quasi-separated, so is its restriction
$f^{-1}(V)\to V$, hence $f^{-1}(V)$ is a quasi-compact and
quasi-separated prescheme (1.2.3); since $Z\cap f^{-1}(V)$ is constructible
in $f^{-1}(V)$ (1.8.1), we see that we may restrict to the case where $Y$ is
affine and $Z$ constructible; $X$ is then quasi-compact and quasi-separated,
hence (1.8.3) there exists a morphism of finite presentation $g:X'\to X$ such
that $Z=g(X')$; since $f\circ g$ is of finite presentation, we see that we
are reduced to the case where $Z=X$; in other words, it will suffice to
prove the
\end{proof}

\begin{lemma}[1.8.4.1]
\label{IV.1.8.4.1}
Let $Y$ be an affine scheme, $f:X\to Y$ a quasi-compact morphism that is
locally of finite presentation; then $f(X)$ is a constructible subset
of $Y$.
\end{lemma}

\begin{proof}
Since $X$ is quasi-compact, it is a finite union of affine open subsets; we
may therefore restrict to the case where $Y=\operatorname{Spec}(A)$,
$X=\operatorname{Spec}(B)$, $B$ being an $A$-algebra of finite presentation
(1.4.6). Now, we have the
\end{proof}

\begin{lemma}[1.8.4.2]
\label{IV.1.8.4.2}
Let $A_0$ be a ring, $(A_\lambda)_{\lambda\in L}$ an inductive system of
$A_0$-algebras, $A=\varinjlim A_\lambda$; if $B$ is an $A$-algebra of finite
presentation, there exist a $\lambda$ and an $A_\lambda$-algebra of finite
presentation $B_\lambda$ such that $B$ is isomorphic to
$B_\lambda\otimes_{A_\lambda}A$.
\end{lemma}

\begin{proof}
By hypothesis, $B$ is isomorphic to an algebra of the form
$A[T_1,\ldots,T_n]/\mathfrak J$, where the $T_i$ are indeterminates and
$\mathfrak J$ an ideal of finite type; let $(F_j)$ ($1\leq j\leq m$) be a
system of generators of $\mathfrak J$. Let
$\varphi_\lambda:A_\lambda\to A$ be the canonical homomorphism. Since $L$ is
filtered, there exists $\lambda$ such that each of the coefficients of each of
the polynomials $F_j$ is the image under $\varphi_\lambda$ of an element of
$A_\lambda$. In other words, there exists a system of polynomials
$(F_{j\lambda})_{1\leq j\leq m}$ of $A_\lambda[T_1,\ldots,T_n]$ such that
$F_j=\varphi_\lambda(F_{j\lambda})$ for $1\leq j\leq m$. If
$\mathfrak J_\lambda$ is the ideal of $A_\lambda[T_1,\ldots,T_n]$ generated
by the $F_{j\lambda}$, $\mathfrak J$ is the image of
$\mathfrak J_\lambda\otimes_{A_\lambda}A$ in
$A[T_1,\ldots,T_n]=A_\lambda[T_1,\ldots,T_n]\otimes_{A_\lambda}A$;
the ring $B_\lambda=A_\lambda[T_1,\ldots,T_n]/\mathfrak J_\lambda$ answers
the question.
\end{proof}

We will apply this lemma here by considering $A$ as the inductive limit of its
sub-$\mathbf Z$-algebras \emph{of finite type}. We therefore see that there
exist such a sub-$\mathbf Z$-algebra $A_0$ of $A$ and an $A_0$-algebra of
finite presentation $B_0$ such that $B$ is isomorphic to
$B_0\otimes_{A_0}A$; set $X_0=\operatorname{Spec}(B_0)$,
$Y_0=\operatorname{Spec}(A_0)$, so that we have $X=X_0\times_{Y_0}Y$,
$f$ being the projection $X\to Y$; let $f_0:X_0\to Y_0$,
$g_0:Y\to Y_0$ be the structural morphisms; it follows from (I, 3.4.8) that
we have $f(X)=g_0^{-1}(f_0(X_0))$; taking into account (1.8.2), it therefore
suffices to show that $f_0(X_0)$ is constructible. In other words, we are
finally reduced to proving the

\begin{corollary}[1.8.5]
\label{IV.1.8.5}
Let $Y$ be a Noetherian prescheme, $f:X\to Y$ a morphism of finite type.
For every constructible subset $Z$ of $X$, $f(Z)$ is a constructible subset
of $Y$.
\end{corollary}

\oldpage[IV-1]{240}

\begin{proof}
We reduce as above to proving that $f(X)$ is constructible.
Let us apply the criterion (0\textsubscript{III}, 9.2.3): we must prove that
for every irreducible closed subset $T$ of $Y$, $T\cap f(X)=f(f^{-1}(T))$ is
rare in $T$ or contains a non-empty open subset of $T$; by taking a
sub-prescheme of $Y$ having $T$ as underlying space (I, 5.2.1) and
considering its inverse image under $f$, we finally see (taking into account
(1.5.4, (iii)) that we are reduced to proving that if we suppose $Y$
\emph{irreducible and $f$ dominant}, $f(X)$ contains a non-empty open subset
of $Y$. We may moreover suppose $Y$ affine; then $X$ is a finite union of
affine open subsets $V_i$, and since $f(X)$ is dense in $Y$, so is at least
one of the $f(V_i)$. We may therefore also suppose $X$ affine; if
$(X_j)$ is the (finite) family of irreducible components of $X$, we see
as above that at least one of the $f(X_j)$ is dense in $Y$; we may therefore
suppose $X$ irreducible. Finally, by replacing $f$ by $f_{\mathrm{red}}$
(1.5.4, (vi)) we may suppose $X$ and $Y$ \emph{integral}. Then
$X=\operatorname{Spec}(B)$, $Y=\operatorname{Spec}(A)$, where $A$ is a
Noetherian integral ring, $B$ an integral $A$-algebra of finite type
containing $A$ (I, 1.2.7). It suffices to show that there exists $g\in A$
such that, for every $y\in D(g)$ (that is, such that $g\notin\mathfrak p_y$)
the ideal $\mathfrak p_y$ is the intersection of $A$ and a prime ideal of $B$,
for this will show that $D(g)\subset f(X)$. Finally it suffices to prove the
\end{proof}

\begin{lemma}[1.8.5.1]
\label{IV.1.8.5.1}
Let $A$ be an integral ring, $B$ an integral $A$-algebra of finite type.
There exists $g\in A$ such that every homomorphism of $A$ into an
algebraically closed field $\Omega$ that is nonzero at $g$ extends to a
homomorphism of $B$ into $\Omega$.
\end{lemma}

\begin{proof}
Now this is a classical result of commutative algebra (Bourbaki,
\emph{Alg. comm.}, chap. V, \S~3, n\textsuperscript{o}~1, cor. 3 of th. 1).
\end{proof}

\begin{corollary}[1.8.6]
\label{IV.1.8.6}
Let $Y$ be an irreducible prescheme, $\eta$ its generic point,
$f:X\to Y$ a morphism locally of finite type. If $f^{-1}(\eta)$ is not
empty, there exists an open neighbourhood $U$ of $\eta$ in $Y$ such that
$f^{-1}(y)$ is non-empty for every $y\in U$. If $\mathscr F$ is a
quasi-coherent $\mathcal O_X$-module of finite type, and if
$\mathscr F_\eta=\mathscr F\otimes_{\mathcal O_Y}k(\eta)$ is not zero,
then there exists an open neighbourhood $U'$ of $\eta$ in $Y$ such that
$\mathscr F_y=\mathscr F\otimes_{\mathcal O_Y}k(y)$ is not zero for
every $y\in U'$.
\end{corollary}

\begin{proof}
If $p_y$ is the canonical projection of the fibre
$f^{-1}(y)=X\times_Y\operatorname{Spec}(k(y))$ into $X$, we have
$\mathscr F_y=p_y^*(\mathscr F)$, hence
$\operatorname{Supp}(\mathscr F_y)=p_y^{-1}(\operatorname{Supp}(\mathscr F))
=\operatorname{Supp}(\mathscr F)\cap f^{-1}(y)$ by virtue of (I, 9.1.13.1) and
(I, 3.6.1), since $\mathscr F$ is of finite type; moreover
$\operatorname{Supp}(\mathscr F)$ is closed in $X$ (0\textsubscript{I},
5.2.2), and if $Z$ is a closed sub-prescheme of $X$ having
$\operatorname{Supp}(\mathscr F)$ as underlying space (I, 5.2.1), and $j$ the
canonical immersion $Z\to X$, $f\circ j$ is locally of finite type (1.3.4);
this proves that the first assertion implies the second. To prove the
first assertion, note first that we may suppose $X$ and $Y$
reduced (1.3.4, (vi)), in other words $Y$ integral. Let $\xi$ be a point of
$f^{-1}(\eta)$; we may replace $X$ and $Y$ by affine open neighbourhoods of
$\xi$ and $\eta$ respectively, and consequently suppose that
$X=\operatorname{Spec}(B)$, $Y=\operatorname{Spec}(A)$, $B$ being an
$A$-algebra of finite type. Moreover, if $Z$ is the reduced closed
sub-prescheme of $X$ having $\overline{\{\xi\}}$ as underlying set
(I, 5.2.1), we may, as above, replace $X$ by $Z$, in other words suppose $B$
\emph{integral} (I, 5.1.4). Since the morphism $f$ is then dominant,
the corresponding homomorphism $\varphi:A\to B$ is injective (I, 1.2.7);
hence the corollary finally follows from Lemma (1.8.5.1).
\end{proof}

\oldpage[IV-1]{241}

\begin{proposition}[1.8.7]
\label{IV.1.8.7}
Let $S$ be a prescheme, $g:X\to S$, $h:Y\to S$ two morphisms of finite
presentation, $f:X\to Y$ an $S$-morphism. For every $s\in S$, set
$X_s=g^{-1}(s)=X\times_S\operatorname{Spec}(k(s))$,
$Y_s=h^{-1}(s)=Y\times_S\operatorname{Spec}(k(s))$,
$f_s=f\times 1:X_s\to Y_s$. Then the set of $s\in S$ such that $f_s$
is surjective (resp. radicial) is locally constructible.
\end{proposition}

\begin{proof}
Let $E$ be the set of $s\in S$ such that $f_s$ is surjective; the set
$S-E$ is equal to $h(Y-f(X))$; now $f$ is of finite presentation
(1.6.2, (v)), hence $f(X)$ is locally constructible in $Y$ (1.8.4); since $h$
is of finite presentation, $h(Y-f(X))$ is locally constructible in $S$, hence
so is $E$.

To show that the set $E'$ of $s\in S$ such that $f_s$ is radicial
is locally constructible, we will use the
\end{proof}

\begin{lemma}[1.8.7.1]\label{IV.1.8.7.1}
\EGAReaderTarget{ega4.c04.lem.1-8-7-1}{IV.1.8.7.1-fr}
Let $U$, $V$ be two preschemes; for a morphism $f:U\to V$ to be
radicial, it is necessary and sufficient that the diagonal morphism
$\Delta_f:U\to U\times_VU$ be surjective. Consequently, every radicial
morphism is separated.
\end{lemma}

\begin{proof}
Indeed, $\Delta_f$, being an immersion (I, 5.3.9), is a morphism
locally of finite type (1.3.4); to say that $\Delta_f$ is surjective
therefore means that for every algebraically closed field $K$, the
corresponding map $U(K)\to(U\times_VU)(K)=U(K)\times_{V(K)}U(K)$ is surjective
(1.3.7 and I, 3.4.2.1); but by the definition of a fibre product of sets,
this means that the map $U(K)\to V(K)$ corresponding to $f$ is
injective, and this is equivalent to saying that $f$ is radicial (I, 3.5.5).

This being so, $\Delta_{f_s}$ is deduced from $\Delta_f:X\to X\times_YX$ by
the base change $\operatorname{Spec}(k(s))\to S$ (I, 5.3.4). Since $f$
is of finite presentation, so is the structural morphism
$X\times_YX\to Y$ (1.6.2, (iv)), hence so is $\Delta_f$ (1.6.2, (v)); it
therefore suffices to apply the first part of the proposition to $\Delta_f$,
using Lemma (1.8.7.1).
\end{proof}


\subsection{Pro-constructible and ind-constructible sets}
\label{subsection:IV.1.9}

\begin{lemma}[1.9.1]
\label{IV.1.9.1}
Let $(A_\alpha)_{\alpha\in I}$ be an inductive system of rings whose index
set $I$ is filtered to the right, and let $A=\varinjlim A_\alpha$. For
$A=0$, it is necessary and sufficient that there exist an index $\gamma$ such
that $A_\gamma=0$ (and we then have $A_\beta=0$ for every $\beta\geq\gamma$).
\end{lemma}

\begin{proof}
Indeed, for every $\alpha\in I$, the canonical homomorphism
$\varphi_\alpha:A_\alpha\to A$ maps the unit element to the unit element;
to say that $A=0$ means that $\varphi_\alpha(1)=0$, or again
$\varphi_\alpha(1)=\varphi_\alpha(0)$; we know that this is equivalent to
saying that there exists an index $\gamma\geq\alpha$ such that the
homomorphism $\varphi_{\gamma\alpha}:A_\alpha\to A_\gamma$ satisfies
$\varphi_{\gamma\alpha}(1)=\varphi_{\gamma\alpha}(0)$, and this last
relation is equivalent to $A_\gamma=0$, whence the lemma.
\end{proof}

\begin{proposition}[1.9.2]
\label{IV.1.9.2}
Let $B$ be a ring, $(A_\alpha)_{\alpha\in I}$ an inductive system of
$B$-algebras whose index set $I$ is filtered to the right, and let
$A=\varinjlim A_\alpha$; set $Y=\operatorname{Spec}(B)$,
$X_\alpha=\operatorname{Spec}(A_\alpha)$, $X=\operatorname{Spec}(A)$, and
let $u:X\to Y$, $u_\alpha:X_\alpha\to Y$ be the morphisms corresponding to
the structural homomorphisms $B\to A$, $B\to A_\alpha$. Then:
\end{proposition}

\begin{env}[1.9.2.i]
\label{IV.1.9.2.i}
\emph{For $X=\varnothing$, it is necessary and sufficient that there exist a
$\gamma\in I$ such that $X_\gamma=\varnothing$, and we then have
$X_\alpha=\varnothing$ for $\alpha\geq\gamma$.}
\end{env}

\oldpage[IV-1]{242}

\begin{env}[1.9.2.ii]
\label{IV.1.9.2.ii}
\emph{We have}
\[
  u(X)=\bigcap_{\alpha\in I}u_\alpha(X_\alpha). \tag{1.9.2.1}
  \label{IV.1.9.2.1}
\]
\end{env}

Assertion (i) is none other than the translation of (1.9.1). To prove
(ii), note that, since $u$ factors as
$X\to X_\alpha\xrightarrow{u_\alpha}Y$, the left-hand side of (1.9.2.1) is
contained in the right-hand side. Conversely, let $y\in Y-u(X)$; we have
$u^{-1}(y)=\varnothing$, and $u^{-1}(y)$ is the underlying space of the
$k(y)$-prescheme $\operatorname{Spec}(A\otimes_Bk(y))$; since in the
category of $B$-modules the functor $\varinjlim$ commutes with the tensor
product, $A\otimes_Bk(y)$ is the inductive limit of the inductive system of
rings $A_\alpha\otimes_Bk(y)$; Lemma (1.9.1) shows that there exists $\alpha$
such that $A_\alpha\otimes_Bk(y)=0$, that is,
$u_\alpha^{-1}(y)=\varnothing$, hence $y\notin u_\alpha(X_\alpha)$. Q.E.D.

\begin{proposition}[1.9.3]
\label{IV.1.9.3}
Let $X$ be a quasi-compact and quasi-separated prescheme, $E$ a subset of
$X$, $(X_\alpha)_{\alpha\in I}$ an affine open covering of $X$, and for
every $\alpha\in I$, set $E_\alpha=E\cap X_\alpha$. The following conditions
are equivalent:
\begin{enumerate}
  \item[a)] There exists a quasi-compact morphism $f:X'\to X$ such that
  $E=f(X')$.
  \item[a$'$)] For every $\alpha$, there exists a quasi-compact morphism
  $f_\alpha:X'_\alpha\to X_\alpha$ such that
  $E_\alpha=f_\alpha(X'_\alpha)$.
  \item[a$''$)] For every $\alpha$, there exist an affine scheme $X''_\alpha$
  and a morphism $f'_\alpha:X''_\alpha\to X_\alpha$ such that
  $f'_\alpha(X''_\alpha)=E_\alpha$.
  \item[b)] $E$ is an intersection of constructible subsets of $X$.
  \item[b$'$)] $F=X-E$ is a union of constructible subsets of $X$.
\end{enumerate}
\end{proposition}

\begin{proof}
It is clear that $b)$ and $b')$ are equivalent. Condition $a)$ implies
$a')$ by taking for $X'_\alpha$ the prescheme induced on the open subset
$f^{-1}(X_\alpha)$ of $X'$ and for $f_\alpha$ the restriction of $f$. To
see that $a')$ implies $a'')$, it suffices to observe that $X'_\alpha$ is a
union of affine open subsets $X''_{\alpha\lambda}$ ($\lambda\in L$); let
$X''_\alpha$ be the \emph{sum} prescheme of the $X''_{\alpha\lambda}$
($\lambda\in L$), which is an affine scheme; if
$g_\alpha:X''_\alpha\to X'_\alpha$ is the morphism coinciding with
the identity on each of the $X''_{\alpha\lambda}$, it is clear that if we
set $f'_\alpha=f_\alpha\circ g_\alpha$, we have
$E_\alpha=f'_\alpha(X''_\alpha)$ since $g_\alpha$ is surjective. To
show that $a'')$ implies $a)$, note that there is a finite subset $J$ of
$I$ such that the $X_\alpha$ with indices $\alpha\in J$ cover $X$; let
$X'$ be the sum prescheme of the $X''_\alpha$ for $\alpha\in J$, and let
$f:X'\to X$ be the morphism coinciding with $j_\alpha\circ f'_\alpha$ for
every $\alpha\in J$, where $j_\alpha:X_\alpha\to X$ is the canonical
injection. Since $E$ is the union of the $E\cap X_\alpha$ for $\alpha\in J$,
we indeed have $E=f(X')$, and it remains to see that $f$ is a quasi-compact
morphism; but the hypothesis that $X$ is quasi-separated implies that
$j_\alpha$ is quasi-compact (1.2.7), hence so is
$j_\alpha\circ f'_\alpha$ (1.1.1 and 1.1.2), and finally so is $f$ (1.1.6).

It therefore remains to prove the equivalence of $a'')$ and $b)$. Let us
first prove that $a'')$ implies $b)$: consider a finite subset $J$ of $I$
such that the $X_\alpha$ for $\alpha\in J$ cover $X$; it will suffice to
show that, for every $\alpha\in J$, $E_\alpha$ is an intersection of
constructible subsets $F_{\alpha\lambda}$ of $X_\alpha$ ($\lambda\in L_\alpha$);
indeed, every $F_{\alpha\lambda}$ is also a constructible subset of $X$
by virtue of (0\textsubscript{III}, 9.1.8, (ii)), since the hypothesis that
$X$ is quasi-separated implies that every $X_\alpha$ is retrocompact in $X$
(1.2.7); $E$, being the union of the $E_\alpha$ for $\alpha\in J$, is the
intersection of the (finite) unions
$\bigcup_{\alpha\in J}F_{\alpha,\lambda(\alpha)}$, for
all choices of the $\lambda(\alpha)\in L_\alpha$; these unions being
constructible in $X$, this establishes our assertion. We may therefore
restrict to the case

\oldpage[IV-1]{243}

where $X=\operatorname{Spec}(A)$ and where $E=f(X')$, where
$X'=\operatorname{Spec}(A')$ is affine, $f:X'\to X$ a morphism
corresponding to a ring homomorphism $A\to A'$. Now note that
$A'$ is the inductive limit of the family $(A'_\lambda)$ of its
sub-$A$-algebras \emph{of finite type}, ordered by inclusion; if
$X'_\lambda=\operatorname{Spec}(A'_\lambda)$, $f$ factors as
$X'\to X'_\lambda\xrightarrow{f_\lambda}X$ and it follows from (1.9.2, (ii))
that $f(X')=\bigcap_\lambda f_\lambda(X'_\lambda)$; if we prove
that $f_\lambda(X'_\lambda)$ is an intersection of constructible subsets of
$X$, the same will hold for $f(X')$. We may therefore restrict to the case
where $A'$ is an $A$-algebra \emph{of finite type}. Now, we have the
following lemma:
\end{proof}

\begin{lemma}[1.9.3.1]
\label{IV.1.9.3.1}
Let $A$ be a ring, $A'$ an $A$-algebra of finite type. Then $A'$ is
$A$-isomorphic to a filtered inductive limit $\varinjlim A''_\mu$ where the
$A''_\mu$ are $A$-algebras of finite presentation (1.4.1).
\end{lemma}

\begin{proof}
Indeed, we can write $A'=B/\mathfrak J$, with
$B=A[T_1,\ldots,T_n]$ and $\mathfrak J$ an ideal of $B$. But $\mathfrak J$
is the inductive limit of the ideals \emph{of finite type} $\mathfrak J'_\mu$
of $B$ contained in $\mathfrak J$, ordered by inclusion; since in the
category of $B$-modules the functor $\varinjlim$ is exact, we conclude
that $A'=\varinjlim B/\mathfrak J'_\mu$ up to an $A$-isomorphism.
\end{proof}

The same argument as above then makes it possible to reduce to the case where
$A'$ is an $A$-algebra of finite presentation; but then $f(X')$ is
constructible in $X$ by virtue of Chevalley's theorem (1.8.5), which
proves $b)$.

Finally, let us prove that $b)$ implies $a'')$. If $E$ is the intersection of
a family $(G_\lambda)_{\lambda\in L}$ of constructible subsets of $X$,
each $E_\alpha$ is the intersection of the $G_\lambda\cap X_\alpha$, which are
constructible in $X_\alpha$ (0\textsubscript{III}, 9.1.8, (i)), hence we
are reduced to the case where $X=\operatorname{Spec}(A)$ is affine.

We then know (1.8.3) that for every $\lambda\in L$ there exists a morphism
$f_\lambda:X'_\lambda\to X$, where
$X'_\lambda=\operatorname{Spec}(A'_\lambda)$ is affine, such that
$f_\lambda(X'_\lambda)=G_\lambda$. It therefore suffices to prove the
following lemma:

\begin{lemma}[1.9.3.2]
\label{IV.1.9.3.2}
Let $(A'_\lambda)_{\lambda\in L}$ be a family of $A$-algebras,
$X=\operatorname{Spec}(A)$, $X'_\lambda=\operatorname{Spec}(A'_\lambda)$,
and let $f_\lambda:X'_\lambda\to X$ be the structural morphism. For every
finite subset $J$ of $L$, set
$A'_J=\bigotimes_{\lambda\in J}A'_\lambda$ (tensor product of
$A$-algebras), $X'_J=\operatorname{Spec}(A'_J)$, and let $f_J:X'_J\to X$ be
the structural morphism. We then have
$f_J(X'_J)=\bigcap_{\lambda\in J}f_\lambda(X'_\lambda)$. If
$A'=\varinjlim A'_J$, $J$ running through the filtered set of finite subsets
of $L$, $X'=\operatorname{Spec}(A')$, and if $f:X'\to X$ is the
structural morphism, we have
$f(X')=\bigcap_{\lambda\in L}f_\lambda(X'_\lambda)$.
\end{lemma}

\begin{proof}
The first assertion follows from (I, 3.4.7); the second follows from (1.9.2,
(ii)), which gives the relation $f(X')=\bigcap_Jf_J(X'_J)$.
\end{proof}

\begin{definition}[1.9.4]
\label{IV.1.9.4}
Let $X$ be a topological space. We say that a subset $E$ of $X$ is
\emph{pro-constructible} (resp. \emph{ind-constructible}) in $X$ if, for
every $x\in X$, there exists an open neighbourhood $U$ of $x$ in $X$ such that
$E\cap U$ is an intersection (resp. a union) of locally constructible subsets
of $U$.
\end{definition}

Taking into account (1.2.7) and (0\textsubscript{III}, 9.1.8 and 9.1.12), the
equivalent conditions of Proposition (1.9.3) can therefore be expressed by
saying that $E$ is \emph{pro-constructible} in $X$, the locally constructible
sets in $X$ being identical with the constructible sets in $X$ under the
hypotheses of (1.9.3).

\oldpage[IV-1]{244}

\begin{proposition}[1.9.5]
\label{IV.1.9.5}
Let $X$ be a prescheme.
\begin{enumerate}
  \item[(i)] For a subset $E$ of $X$ to be pro-constructible, it is necessary
  and sufficient that $X-E$ be ind-constructible.
  \item[(ii)] Every finite union and every intersection of
  pro-constructible subsets of $X$ is pro-constructible. Every finite
  intersection and every union of ind-constructible subsets of $X$ is
  ind-constructible.
  \item[(iii)] Every intersection (resp. union) of locally constructible
  subsets of $X$ is pro-constructible (resp. ind-constructible).
  Conversely, if $X$ is quasi-compact and quasi-separated, every
  pro-constructible (resp. ind-constructible) subset of $X$ is an
  intersection (resp. a union) of constructible subsets of $X$.
  \item[(iv)] Let $(U_\alpha)$ be an open covering of $X$. For a subset $E$
  of $X$ to be pro-constructible (resp. ind-constructible) in $X$, it is
  necessary and sufficient that for every $\alpha$, $E\cap U_\alpha$ be
  pro-constructible (resp. ind-constructible) in $U_\alpha$.
  \item[(v)] Every pro-constructible subset $E$ of $X$ is retrocompact
  in $X$. For a locally closed subset $E$ of $X$ to be retrocompact in $X$,
  it is necessary and sufficient that it be pro-constructible in $X$.
  \item[(vi)] Let $f:X'\to X$ be a morphism of preschemes; for every
  pro-constructible (resp. ind-constructible) subset $E$ of $X$, $f^{-1}(E)$
  is pro-constructible (resp. ind-constructible) in $X'$.
  \item[(vii)] Let $f:X'\to X$ be a quasi-compact morphism; for every
  pro-constructible subset $E'$ of $X'$, $f(E')$ is pro-constructible in
  $X$; in particular $f(X')$ is pro-constructible in $X$.
  \item[(viii)] Let $f:X'\to X$ be a morphism locally of finite
  presentation; for every ind-constructible subset $E'$ of $X'$, $f(E')$ is
  ind-constructible in $X$; in particular $f(X')$ is ind-constructible
  in $X$.
  \item[(ix)] Suppose that $X$ is quasi-compact and quasi-separated; then,
  for a subset $E$ of $X$ to be pro-constructible, it is necessary and
  sufficient that there exist an affine scheme $X'$ and a (necessarily
  quasi-compact) morphism $f:X'\to X$ such that $E=f(X')$.
\end{enumerate}
\end{proposition}

\begin{proof}
Properties (i), (ii), (iv) and the first assertion of (iii) follow
from Definition (1.9.4) and are valid for an arbitrary topological space,
by using the mutual distributivity of intersection and union and the fact
that if $T$ is locally constructible in $X$, $T\cap U$ is locally
constructible in $U$ for every open subset $U$ of $X$.
If $X$ is quasi-compact and quasi-separated and $E$ is pro-constructible in
$X$, there is a finite covering $(U_i)$ of $X$ consisting of affine open
subsets such that $E\cap U_i$ is an intersection of constructible subsets
$M_{i\lambda}$ of $U_i$ ($\lambda\in L_i$); the $M_{i\lambda}$ are also
constructible in $X$ by virtue of (1.2.7) and (0\textsubscript{III}, 9.1.8,
(ii)), and $E$ is the intersection of the finite unions
$\bigcup_iM_{i,\lambda(i)}$ (where $\lambda(i)\in L_i$ for every $i$), which
are constructible subsets of $X$; this proves the second assertion of (iii).
Assertion (ix) follows from (iii) and from (1.9.3). To prove the first
assertion of (v), we may restrict to the case where $X$ is affine, and then
prove that $E$ is quasi-compact (0\textsubscript{III}, 9.1.1); but since $X$
is then quasi-separated, there exists by virtue of (ix) a quasi-compact
morphism $f:X'\to X$ such that $E=f(X')$; since $X'$ is quasi-compact, so is
its image $E$ under a continuous map.

To prove (vi), we may restrict, by virtue of (iv), to the case where $X$ is
affine; the conclusion then follows from (iii) and from (1.8.2).

\oldpage[IV-1]{245}

Likewise, to prove (vii), we may restrict to the case where $X$ is affine;
then $X'$ is a finite union of affine open subsets $X'_i$ and $f(E')$ is the
union of the $f(E'\cap X'_i)$, so that we may also suppose that $X'$ is
affine, by virtue of (ii); but then $E'=g(X'')$, where $g:X''\to X'$ is a
quasi-compact morphism, by virtue of (ix), and we have $f(E')=f(g(X''))$,
hence $f(E')$ is pro-constructible since $g\circ f$ is quasi-compact.

We deduce from (vii) the proof of the second assertion of (v): indeed,
let $E$ be a locally closed and retrocompact subset of $X$, and consider a
sub-prescheme of $X$ having $E$ as underlying space (I,
5.2.1); the canonical injection $j:E\to X$ is then by hypothesis a
quasi-compact morphism, and it follows from (vii) that $E=j(E)$ is
pro-constructible in $X$.

Finally, to prove (viii), we may again suppose $X$ affine; moreover, if
$(X'_\alpha)$ is an affine open covering of $X'$ such that
$E'\cap X'_\alpha$ is a union of constructible subsets of $X'_\alpha$
((iii) and (iv)), $f(E')$ is the union of the $f(E'\cap X'_\alpha)$ and, by
virtue of Chevalley's theorem (1.8.4), each of the $f(E'\cap X'_\alpha)$ is
a union of constructible subsets of $X$, whence the conclusion.
\end{proof}

\begin{env}[1.9.6]
\label{IV.1.9.6}
\textbf{Examples} (1.9.6). --- Every \emph{finite} subset of a prescheme $X$
is pro-constructible: indeed, it suffices to consider the subsets
$\{x\}$ consisting of a single point (1.9.5, (ii)), and $\{x\}$ is the image
of $\operatorname{Spec}(k(x))$ under the canonical morphism
$\operatorname{Spec}(k(x))\to X$, which is quasi-compact; whence the
conclusion by (1.9.5, (vii)). On the other hand, a subset $\{x\}$ consisting
of one point is not necessarily ind-constructible; for example, let $A$ be a
Noetherian integral ring having infinitely many maximal ideals, and let $x$
be the generic point of $X=\operatorname{Spec}(A)$; if $\{x\}$ were
ind-constructible in $X$, it would be constructible by virtue of (1.9.5,
(iii)) and would consequently contain a non-empty open subset of $X$
(0\textsubscript{III}, 9.2.2); but this contradicts the hypothesis that $A$
has infinitely many maximal ideals, by virtue of the Artin--Tate theorem
(0, 16.3.3).

Every \emph{closed} subset $Y$ of a prescheme $X$ is pro-constructible, by
(1.9.5, (v)). Every open subset of $X$ is therefore ind-constructible, but
an open subset $U$ of $X$ is pro-constructible only if it is
\emph{retrocompact}, again by virtue of (1.9.5, (v)).

Finally, if $X$ is a \emph{Noetherian} prescheme, hence quasi-separated
(1.2.8), it follows from (1.9.5, (iii)) and from (0\textsubscript{III}, 9.1.7)
that the ind-constructible subsets of $X$ are the \emph{unions of locally
closed subsets of $X$}.
\end{env}

\begin{theorem}[1.9.7]
\label{IV.1.9.7}
Let $X$ be a quasi-compact prescheme, $E$ an ind-constructible subset of
$X$, $(F_\lambda)_{\lambda\in L}$ a family of pro-constructible subsets
of $X$ such that $\bigcap_{\lambda\in L}F_\lambda\subset E$; then there
exists a finite subset $J$ of $L$ such that
$\bigcap_{\lambda\in J}F_\lambda\subset E$.
\end{theorem}

\begin{proof}
Note that the sets $F=X-E$ and $F\cap F_\lambda$ are
pro-constructible, hence we are reduced to the
\end{proof}

\begin{corollary}[1.9.8]
\label{IV.1.9.8}
Let $X$ be a quasi-compact prescheme, $(F_\lambda)_{\lambda\in L}$ a
family of pro-constructible subsets of $X$ whose intersection is empty;
then there exists a finite subfamily $(F_\lambda)_{\lambda\in J}$ whose
intersection is empty.
\end{corollary}

\begin{proof}
Covering $X$ by finitely many affine open subsets, and using (1.9.5,
(iv)), we are reduced to the case where $X$ is affine; by virtue of (1.9.5,
(ix)), we have $F_\lambda=f_\lambda(X'_\lambda)$,

\oldpage[IV-1]{246}

where $X'_\lambda=\operatorname{Spec}(A'_\lambda)$ is affine, and $f_\lambda$
is a morphism $X'_\lambda\to X$; then, with the notation of (1.9.3.2),
we have by hypothesis $f(X')=\varnothing$, hence $X'=\varnothing$; but this
implies by (1.9.2, (i)) that there exists a finite subset $J$ of $L$ such
that $X'_J=\varnothing$, hence
$f_J(X'_J)=\bigcap_{\lambda\in J}F_\lambda=\varnothing$.
\end{proof}

\begin{corollary}[1.9.9]
\label{IV.1.9.9}
Let $X$ be a quasi-compact prescheme, $F$ a pro-constructible subset of
$X$, $(E_\lambda)_{\lambda\in L}$ a family of ind-constructible subsets
of $X$ such that $F\subset\bigcup_{\lambda\in L}E_\lambda$; then there
exists a finite subset $J$ of $L$ such that
$F\subset\bigcup_{\lambda\in J}E_\lambda$. In particular, from every
covering of $X$ by ind-constructible subsets, we can extract a finite
covering of $X$.
\end{corollary}

\begin{proof}
This follows from (1.9.7) by passing to complements.
\end{proof}

\begin{proposition}[1.9.10]
\label{IV.1.9.10}
Let $X$ be a prescheme. For a subset $E$ of $X$ to be
ind-constructible, it is necessary that, for every $x\in X$, the intersection
$E\cap\overline{\{x\}}$ be a neighbourhood of $x$ in
$\overline{\{x\}}$. If $X$ is locally Noetherian, this condition is
sufficient.
\end{proposition}

\begin{proof}
Suppose $E$ ind-constructible, and let $Y$ be the intersection of
$\overline{\{x\}}$ and an affine open subset of $X$ containing $x$; there is
therefore a sub-prescheme of $X$ having $Y$ as underlying space (I, 5.2.1),
and if $j:Y\to X$ is the canonical injection,
$E\cap Y=j^{-1}(E)$ is ind-constructible in $Y$ (1.9.5, (vi)). Since $Y$
is quasi-compact and separated, $E\cap Y$ is therefore a union of
constructible subsets of $Y$ (1.9.5, (iii)); consequently, there are two
open subsets $U$, $V$ of $Y$ such that
$x\in U\cap\complement V\subset E\cap Y$ (0\textsubscript{III}, 9.1.3).
But since $x$ is the generic point of the irreducible space $Y$, $V$ is
necessarily empty, and we have $U\subset E\cap Y$.

Now suppose $X$ locally Noetherian, and let $E$ be a subset of
$X$ satisfying the condition of the statement. By virtue of Definition
(1.9.4), we may restrict to the case where $X$ is Noetherian. Then, for every
$x\in E$, $E\cap\overline{\{x\}}$ contains a non-empty subset of the form
$U\cap\overline{\{x\}}$, where $U$ is open in $X$; this shows that $E$
is a union of locally closed subsets of $X$, hence is
ind-constructible (1.9.6).
\end{proof}

\begin{proposition}[1.9.11]
\label{IV.1.9.11}
Let $X$ be a prescheme, $E$ a subset of $X$. The following properties are
equivalent:
\begin{enumerate}
  \item[a)] $E$ is locally constructible.
  \item[b)] $E$ is ind-constructible and pro-constructible.
  \item[c)] $E$ and $X-E$ are pro-constructible.
  \item[d)] $E$ and $X-E$ are ind-constructible.
\end{enumerate}
\end{proposition}

\begin{proof}
It evidently suffices to prove that $d)$ implies $a)$, and we may restrict
to the case where $X$ is affine. Then (1.9.5, (iii)), $E$ (resp. $X-E$) is
the union of a family $(E_\lambda)$ (resp. $(E'_\mu)$) of constructible
subsets of $X$; since the $E_\lambda$ and the $E'_\mu$ form a covering of
$X$, it follows from (1.9.9) that there are finitely many indices
$\lambda_i$, $\mu_j$ such that the $E_{\lambda_i}$ and the $E'_{\mu_j}$ form
a covering of $X$; this implies that $E$ is the union of the
$E_{\lambda_i}$, hence is constructible.
\end{proof}

\begin{corollary}[1.9.12]
\label{IV.1.9.12}
Let $f:X\to Y$ be a surjective morphism which is either quasi-compact or
locally of finite presentation. For a subset $E$ of $Y$ to be
locally constructible (resp. pro-constructible, resp. ind-constructible),
it is necessary and sufficient that $f^{-1}(E)$ be so in $X$.
\end{corollary}

\oldpage[IV-1]{247}

\begin{proof}
We know that the condition is necessary ((1.8.2) and (1.9.5, (vi))); to
show that it is sufficient, we are reduced to the case where $X$ is affine,
by virtue of (1.9.5, (iv)); moreover, by virtue of (1.9.11), we may restrict
to the case where $f^{-1}(E)$ is pro-constructible, \emph{or} to the case
where $f^{-1}(E)$ is ind-constructible. If $f$ is surjective and
quasi-compact, and $f^{-1}(E)$ pro-constructible, it follows from (1.9.5,
(vii)) that $E=f(f^{-1}(E))$ is pro-constructible. If $f$ is surjective and
locally of finite presentation, and $f^{-1}(E)$ ind-constructible,
$E=f(f^{-1}(E))$ is ind-constructible by (1.9.5, (viii)), which completes the
proof.
\end{proof}

\begin{env}[1.9.13]
\label{IV.1.9.13}
Let $X$ be a prescheme, $\mathscr T$ its topology; it follows from (1.9.5,
(i) and (ii)) that the ind-constructible (resp. pro-constructible) subsets of
$X$ are the open (resp. closed) subsets for a topology on $X$,
called the \emph{constructible topology}, which we will denote by
$\mathscr T^{\mathrm{cons}}$; we will denote by $X^{\mathrm{cons}}$
the set $X$ endowed with the topology $\mathscr T^{\mathrm{cons}}$.
\end{env}

\begin{proposition}[1.9.14]
\label{IV.1.9.14}
Let $X$ be a prescheme, $\mathscr T$ its topology,
$\mathscr T^{\mathrm{cons}}$ the constructible topology on $X$.
\begin{enumerate}
  \item[(i)] The topology $\mathscr T^{\mathrm{cons}}$ is finer than
  $\mathscr T$.
  \item[(ii)] The locally constructible subsets of $X$ are identical
  with the subsets of the space $X^{\mathrm{cons}}$ that are both open and
  closed.
  \item[(iii)] For every morphism $f:X\to Y$, the underlying map from
  $X^{\mathrm{cons}}$ to $Y^{\mathrm{cons}}$ is continuous; we denote it by
  $f^{\mathrm{cons}}$.
  \item[(iv)] If the morphism $f:X\to Y$ is quasi-compact,
  $f^{\mathrm{cons}}$ is a closed map; in particular, if $f$
  is quasi-compact and bijective, $f^{\mathrm{cons}}$ is a homeomorphism.
  \item[(v)] If a morphism $f:X\to Y$ is locally of finite
  presentation, $f^{\mathrm{cons}}$ is an open map; in particular,
  if $f$ is bijective and locally of finite presentation,
  $f^{\mathrm{cons}}$ is a homeomorphism.
  \item[(vi)] For every open subset $U$ of $X$, the topology induced by
  $\mathscr T^{\mathrm{cons}}$ on $U$ is identical with the topology of
  $U^{\mathrm{cons}}$.
\end{enumerate}
\end{proposition}

\begin{proof}
Indeed, (i) follows from the fact that every open subset for $\mathscr T$ is
an open subset for $\mathscr T^{\mathrm{cons}}$ (1.9.6), and (ii) is a
translation of (1.9.11); (iii), (iv) and (v) translate respectively
(1.9.5), (vi), (vii) and (viii). Finally, to prove (vi), it suffices to
observe that the canonical injection $j:U\to X$ is locally of finite
presentation (1.4.3), hence
$j^{\mathrm{cons}}:U^{\mathrm{cons}}\to X^{\mathrm{cons}}$ is a
continuous open map.
\end{proof}

\begin{proposition}[1.9.15]
\label{IV.1.9.15}
Let $X$ be a prescheme.
\begin{enumerate}
  \item[(i)] For $X^{\mathrm{cons}}$ to be quasi-compact, it is necessary and
  sufficient that $X$ be quasi-compact.
  \item[(ii)] For $X^{\mathrm{cons}}$ to be separated, it is necessary and
  sufficient that $X$ be quasi-separated, and $X^{\mathrm{cons}}$ is then
  locally compact and totally disconnected.
  \item[(iii)] For $X^{\mathrm{cons}}$ to be compact, it is necessary and
  sufficient that $X$ be quasi-compact and quasi-separated.
  \item[(iv)] Every point of $X$ has, for the topology
  $\mathscr T^{\mathrm{cons}}$, an open and compact neighbourhood.
  \item[(v)] For a morphism $f:X\to Y$ to be quasi-compact, it is necessary
  and sufficient that the continuous map
  $f^{\mathrm{cons}}:X^{\mathrm{cons}}\to Y^{\mathrm{cons}}$ be proper.
\end{enumerate}
\end{proposition}

\begin{proof}
(i) Since $\mathscr T^{\mathrm{cons}}$ is finer than $\mathscr T$, it is
clear that if $X^{\mathrm{cons}}$ is quasi-compact, so is $X$;
the converse follows from (1.9.9).

(ii) Suppose $X$ quasi-separated, and let us show that $X^{\mathrm{cons}}$ is
totally disconnected: indeed, if $x$, $y$ are two distinct points of
$X$ and if for example $x\notin\overline{\{y\}}$, there exists an
\oldpage[IV-1]{248}
affine open neighbourhood $U$ of $X$ not containing $y$; since $X$ is
quasi-separated, $U$ is retrocompact in $X$ (1.2.7), hence $U$ and $X-U$ are
locally constructible in $X$, and consequently open in
$X^{\mathrm{cons}}$ by virtue of (1.9.11), whence our assertion. Since on
every affine open subset $U$ of $X$ the topology induced by
$\mathscr T^{\mathrm{cons}}$ is that of $U^{\mathrm{cons}}$ by virtue of
(1.9.14, (vi)), it follows from what precedes that $X^{\mathrm{cons}}$ is
locally compact, since $U$ is open in $X^{\mathrm{cons}}$ and
$U^{\mathrm{cons}}$ is compact. It remains to prove that if
$X^{\mathrm{cons}}$ is separated, then $X$ is quasi-separated; indeed,
consider an affine open subset $U$ of $X$; the topology induced on
$U$ by $\mathscr T^{\mathrm{cons}}$ is that of $U^{\mathrm{cons}}$
(1.9.14, (vi)), hence $U$ is compact for this induced topology, since
$U^{\mathrm{cons}}$ is compact by the first part of the argument. If
$V$ is a second affine open subset of $X$, $U\cap V$ is therefore a compact
subset of the separated space $X^{\mathrm{cons}}$, being the intersection of
two compact subsets of this space; since the topology induced by
$\mathscr T^{\mathrm{cons}}$ on $U\cap V$ is that of
$(U\cap V)^{\mathrm{cons}}$ (1.9.14, (vi)) and the identity map
$(U\cap V)^{\mathrm{cons}}\to U\cap V$ is continuous, we conclude that
$U\cap V$ is quasi-compact for the topology induced by $\mathscr T$; $X$
is consequently quasi-separated by virtue of (1.2.7).

(iii) is an immediate corollary of (i) and (ii).

(iv) For every $x\in X$, an affine open neighbourhood $U$ of $x$ for
$\mathscr T$ is also a neighbourhood of $x$ for $\mathscr T^{\mathrm{cons}}$,
and it is compact by virtue of (iii) and of (1.9.14, (vi)), the topology
induced on $U$ by $\mathscr T^{\mathrm{cons}}$ being identical with that of
$U^{\mathrm{cons}}$.

(v) Suppose $f$ quasi-compact; then we already know (1.9.14, (iv)) that
$f^{\mathrm{cons}}$ is a closed map. On the other hand, let $y$ be a
point of $Y$, and set
\[
  Z=f^{-1}(y)=X\times_Y\operatorname{Spec}(k(y));
\]
$Z$ is quasi-compact (1.1.2, (iii)), and since the canonical morphism
$p:Z\to X$ is injective, it follows from (i) and from the fact that the map
$p^{\mathrm{cons}}$ is continuous that the topology induced on $f^{-1}(y)$
by that of $X^{\mathrm{cons}}$ makes $f^{-1}(y)$ a quasi-compact space.
This proves that $f^{\mathrm{cons}}$ is a proper map (Bourbaki,
\emph{Top. gén.}, chap. I\textsuperscript{er}, 3\textsuperscript{e} éd.,
\S~10, n\textsuperscript{o}~2, th.~1). Conversely, suppose that
the continuous map $f^{\mathrm{cons}}$ is proper, and let $V$ be a
quasi-compact open subset of $Y$; if $h:V\to Y$ is the canonical injection,
$h^{\mathrm{cons}}:V^{\mathrm{cons}}\to Y^{\mathrm{cons}}$ is continuous and
injective and $V^{\mathrm{cons}}$ is quasi-compact by (i), hence the topology
induced on $V$ by that of $Y^{\mathrm{cons}}$ makes $V$ a
quasi-compact space. The hypothesis that $f^{\mathrm{cons}}$ is proper then
implies that the topology induced on $f^{-1}(V)$ by that of
$X^{\mathrm{cons}}$ makes $f^{-1}(V)$ a quasi-compact space
(\emph{loc. cit.}, prop.~6), hence $f^{-1}(V)$ is also a quasi-compact
subspace of $X$, which shows that the morphism $f$ is quasi-compact.
\end{proof}

\begin{env}[1.9.16]
\label{IV.1.9.16}
We will show later how, for every prescheme $X$, one can endow the space
$X^{\mathrm{cons}}$ with a sheaf of rings that makes it a prescheme
whose local rings are all fields, identical with the residue fields at the
points of $X$. Such preschemes arise, for example, in a natural way in the
translation, into the language of schemes, of the constructions of
Néron [32] in his theory of the reduction of abelian varieties.
\end{env}

\oldpage[IV-1]{249}

\subsection{Application to open morphisms}
\label{subsection:IV.1.10}

\begin{theorem}[1.10.1]
\label{IV.1.10.1}
Let $X$ be a prescheme, $E$ an ind-constructible subset of $X$ (1.9.4), $x$
a point of $X$. For $x$ to be interior to $E$, it is necessary and
sufficient that every generization ($0_{\mathrm I}$, 2.1.2) $x'$ of $x$
belong to $E$, in other words (I, 2.4.2) that we have
$\operatorname{Spec}(\mathcal O_{X,x})\subset E$.
\end{theorem}

\begin{proof}
The condition is evidently necessary, every neighbourhood of $x$ containing
the generizations of $x$. To see that it is sufficient, we may evidently
restrict to the case where $X=\operatorname{Spec}(A)$ is affine, $x$ being a prime ideal
$\mathfrak p$ of $A$. We then know (1.9.5, (ix)) that there exist an affine
scheme $X'=\operatorname{Spec}(A')$ and a morphism $f:X'\to X$ such that
$f(X')=X-E$. Set $Y=\operatorname{Spec}(\mathcal O_{X,x})$ and
$Y'=X'\times_XY$; the hypothesis means (by
virtue of (I, 3.6.5)) that $Y'=\varnothing$, in other words
$A'\otimes_AA_{\mathfrak p}=0$. Since
$A_{\mathfrak p}=\varinjlim A_t$, where $t$ runs through $A-\mathfrak p$
($0_{\mathrm I}$, 1.4.5), we have $\varinjlim A_t'=0$, the functor
$\varinjlim$ commuting with the tensor product. Consequently (1.9.1), there
exists a $t\in A-\mathfrak p$ such that $A_t'=0$, and $\mathrm D(t)$ is then
an open neighbourhood of $x$ in $X$ contained in $E$.
\end{proof}

\begin{corollary}[1.10.2]
\label{IV.1.10.2}
Let $f:X\to Y$ be a morphism locally of finite presentation, and let
$y\in Y$. For $y$ to be interior to $f(X)$ it is necessary and sufficient that
every generization $y'$ of $y$ belong to $f(X)$.
\end{corollary}

\begin{proof}
Indeed, we know (1.9.5, (viii)) that $f(X)$ is ind-constructible in $Y$, and
it suffices to apply (1.10.1).
\end{proof}

We will say that a continuous map $f:X\to Y$ is \emph{open at a
point} $x\in X$ if the image under $f$ of \emph{every} neighbourhood of $x$
in $X$ is a neighbourhood of $f(x)$ in $Y$.

\begin{proposition}[1.10.3]\label{IV.1.10.3}
\EGAReaderTarget{ega4.c04.prop.1-10-3}{IV.1.10.3-fr}
Let $f:X\to Y$ be a morphism, $x$ a point of $X$, $y=f(x)$. Consider
the following conditions:
\begin{enumerate}
  \item[a)] $f$ is open at the point $x$.
  \item[b)] For every generization $y'$ of $y$, there exists a generization
  $x'$ of $x$ such that $f(x')=y'$.
  \item[b')] The image under $f$ of $\operatorname{Spec}(\mathcal O_{X,x})$ is
  $\operatorname{Spec}(\mathcal O_{Y,y})$.
  \item[c)] For every irreducible closed subset $Y'$ of $Y$ containing $y$,
  there exists an irreducible component of $X'=f^{-1}(Y')$ containing $x$ and
  dominating $Y'$.
\end{enumerate}
Then we have the implications
$\mathrm{a)}\Rightarrow\mathrm{b)}\Leftrightarrow\mathrm{b')}
\Leftrightarrow\mathrm{c)}$. If moreover $f$ is locally of finite
presentation, the four conditions are equivalent.
\end{proposition}

\begin{proof}
By the definition of a generization, the image under $f$ of
$\operatorname{Spec}(\mathcal O_{X,x})$ is always contained in
$\operatorname{Spec}(\mathcal O_{Y,y})$, hence b) and b') are equivalent. It
is immediate
that c) implies b), for if $y'$ is a generization of $y$,
$Y'=\overline{\{y'\}}$ is an irreducible subset of $Y$ containing $y$; if
$x'$ is the generic point of an irreducible component of $f^{-1}(Y')$ that
contains $x$ and dominates $Y'$, we have $f(x')=y'$ ($0_{\mathrm I}$, 2.1.5)
and $x'$ is a generization of $x$. Conversely, b) implies c): indeed, let
$y'$ be the generic point of $Y'$ and let $x'$ be a generization of $x$ such
that $f(x')=y'$; let $Z'$ be an irreducible component of $f^{-1}(Y')$
containing $x'$ (hence $x$); since $f(x')$ is already dense in $Y'$, this is
\emph{a fortiori} the case for $f(Z')$.

To see that a) implies b'), we may evidently restrict to the case where
$X=\operatorname{Spec}(B)$ and $Y=\operatorname{Spec}(A)$ are affine, $x$
being a prime ideal
$\mathfrak q$ of $B$, so that
$\mathcal O_{X,x}=B_{\mathfrak q}=\varinjlim B_t$,
\oldpage[IV-1]{250}
where $t$ runs through $B-\mathfrak q$. Consequently (1.9.2), we have
\[
  f(\operatorname{Spec}(\mathcal O_{X,x}))
  =\bigcap_t f(\operatorname{Spec}(B_t))=\bigcap_t f(\mathrm D(t)),
\]
and by hypothesis, for every $t\in B-\mathfrak q$, $y$ is interior to
$f(\mathrm D(t))$, hence $f(\mathrm D(t))$ contains
$\operatorname{Spec}(\mathcal O_{Y,y})$; whence
$\operatorname{Spec}(\mathcal O_{Y,y})
\subset f(\operatorname{Spec}(\mathcal O_{X,x}))$, which establishes
our assertion. Finally, when $f$ is locally of finite presentation, b)
implies a) by virtue of (1.10.2) applied to the restriction $U\to Y$ of $f$
to an arbitrary open neighbourhood $U$ of $x$.
\end{proof}

\begin{corollary}[1.10.4]
\label{IV.1.10.4}
Let $f:X\to Y$ be a morphism. Consider the following conditions:
\begin{enumerate}
  \item[a)] $f$ is open.
  \item[b)] For every $x\in X$ and every generization $y'$ of $y=f(x)$, there
  exists a generization $x'$ of $x$ such that $f(x')=y'$.
  \item[b')] For every $x\in X$, the image under $f$ of
  $\operatorname{Spec}(\mathcal O_{X,x})$ is
  $\operatorname{Spec}(\mathcal O_{Y,f(x)})$.
  \item[c)] For every irreducible closed subset $Y'$ of $Y$, every
  irreducible component of $f^{-1}(Y')$ dominates $Y'$.
\end{enumerate}
Then we have the implications
$\mathrm{a)}\Rightarrow\mathrm{b)}\Leftrightarrow\mathrm{b')}
\Leftrightarrow\mathrm{c)}$. If moreover $f$ is locally of finite
presentation, the four conditions are equivalent.
\end{corollary}

\begin{proof}
To say that $f$ is open means that $f$ is open at every point $x\in X$,
hence the implications
$\mathrm{a)}\Rightarrow\mathrm{b)}\Leftrightarrow\mathrm{b')}$ follow from
the analogous implications in (1.10.3), as does the implication
$\mathrm{b)}\Rightarrow\mathrm{a)}$ when $f$ is locally of finite
presentation. The implication $\mathrm{c)}\Rightarrow\mathrm{b)}$ also
follows from the analogous implication in (1.10.3); let us finally prove that
b) implies c). Indeed, let $x'$ be the generic point of an irreducible
component of $f^{-1}(Y')$, and let us show that $y'=f(x')$ is the generic
point of $Y'$. Let $y''$ be a generization of $y'$; by hypothesis there
exists a generization $x''$ of $x'$ such that $f(x'')=y''$, and since
$x''\in f^{-1}(Y')$, we necessarily have $x''=x'$, hence $y''=y'$, which
completes the proof.
\end{proof}

\begin{center}
\emph{(To be continued.)}
\end{center}

\clearpage
\oldpage[IV-1]{251}
\phantomsection
\addcontentsline{toc}{section}{Bibliography (continued)}
